\documentclass[11pt]{article}
\usepackage[T1]{fontenc}
\usepackage[utf8]{inputenc}
\usepackage{amsmath,amsthm,mathtools}
\usepackage{newtxtext,newtxmath}
\usepackage[margin=1in]{geometry}
\usepackage{microtype,booktabs,array,longtable}
\usepackage{enumitem}
\usepackage{placeins,needspace}
\usepackage{xurl}
\usepackage[hidelinks,bookmarksnumbered]{hyperref}
\usepackage{fancyhdr}
\usepackage{listings}
\setlength{\headheight}{14pt}
\pagestyle{fancy}
\fancyhf{}
\fancyhead[L]{\small Rational channels beyond finite tracial baths}
\fancyhead[R]{\small Version 1.0}
\fancyfoot[C]{\thepage}
\renewcommand{\headrulewidth}{0.3pt}
\linespread{1.04}
\setlength{\parskip}{3pt}
\setlength{\parindent}{1.2em}
\setlist{itemsep=2pt,topsep=4pt}
\lstset{basicstyle=\small\ttfamily,breaklines=true,columns=fullflexible,keepspaces=true,
  frame=single,framerule=0.3pt,xleftmargin=0.4em,xrightmargin=0.4em}
\newtheorem{theorem}{Theorem}
\newtheorem{proposition}[theorem]{Proposition}
\newtheorem{lemma}[theorem]{Lemma}
\newtheorem{corollary}[theorem]{Corollary}
\theoremstyle{definition}
\newtheorem{definition}[theorem]{Definition}
\newtheorem{remark}[theorem]{Remark}
\newcommand{\C}{\mathbb C}
\newcommand{\Q}{\mathbb Q}
\newcommand{\Z}{\mathbb Z}
\newcommand{\F}{\mathbb F}
\newcommand{\N}{\mathbb N}
\newcommand{\Dhat}{\widehat D}
\newcommand{\Ghat}{\widehat G}
\newcommand{\Hhat}{\widehat H}
\newcommand{\M}{\mathcal M}
\newcommand{\Lab}{\mathcal L}
\newcommand{\Rel}{\mathcal R}
\newcommand{\Tri}{\mathcal T}
\newcommand{\Fact}{\operatorname{Fact}}
\newcommand{\Tr}{\operatorname{Tr}}
\newcommand{\id}{\operatorname{id}}
\newcommand{\supp}{\operatorname{supp}}
\newcommand{\rank}{\operatorname{rank}}
\newcommand{\dist}{\operatorname{dist}}
\newcommand{\Ad}{\operatorname{Ad}}
\newcommand{\SL}{\operatorname{SL}}
\newcommand{\St}{\operatorname{St}}
\newcommand{\U}{\operatorname{U}}
\newcommand{\nnorm}[1]{\left\lVert #1\right\rVert}
\newcommand{\abs}[1]{\left\lvert #1\right\rvert}
\newcommand{\ket}[1]{\lvert #1\rangle}
\newcommand{\bra}[1]{\langle #1\rvert}
\newcommand{\vket}[1]{\lvert #1\rangle\!\rangle}
\newcommand{\vbra}[1]{\langle\!\langle #1\rvert}
\newcommand{\full}{\mathrm{ref}}
\newcommand{\dJ}{d_J}
\newcommand{\dd}{d_\diamond}
\hypersetup{pdfauthor={Nidhal Mghirbi \& Seth Douglas}}

\hypersetup{pdftitle={A Rational Quantum Channel Beyond Finite Tracial Baths: a construction on fifteen system qubits},
pdfauthor={Nidhal Mghirbi and Seth Douglas}}
\title{A Rational Quantum Channel Beyond Finite Tracial Baths\\[0.3em]
\large A construction on fifteen system qubits}
\author{Nidhal Mghirbi \and Seth Douglas}
\date{October 2026 (version 1.0)}
\begin{document}
\maketitle
\begin{abstract}
A quantum channel on a finite system always has a finite realization with a pure environment.
With a maximally mixed environment this can fail even approximately: since
$\mathrm{MIP}^*=\mathrm{RE}$, some factorizable channels lie outside the closure of the finite
matrix-tracial factorizations. We construct an explicit one: a unital channel on fifteen qubits
with 11\,230 Kraus operators, with entries in $\{0,1,4/5,\pm3/5\}$, that factorizes exactly
through a type $\mathrm{II}_1$ factor. The construction starts from an explicit finite
presentation, with 43 generators and 4\,417 relators, of a group that maps into the
amalgamated double of a Kun--Thom pair. In this presentation Thom's witness commutator is a
nontrivial word of length six, and the recent centralizer-rigidity theorems of Alekseev, Liu
and Thom and of Thom imply that every tracial matrix-ultraproduct representation kills it.
Rational two-level blocks enforce the relations, and the Choi matrix gives the word a trace the
identity cannot have. We also give a support-restricted separation of $1/2$ in half-diamond
distance. For a fully referenced version of the channel, any uniform relator tolerance, whose
value is still open, converts into explicit global distance bounds and into a linear lower
bound on the memory needed for repeated use at small error and bounded exchange. Exact data and
independent checks accompany the proofs.
\end{abstract}

\section{The question and the construction}

With a freely chosen initial environment, every channel on a finite-dimensional system has an
exact finite realization~\cite{stinespring}. Requiring the environment to start maximally
mixed changes the question: the bath may be as large as we like, but it contributes no purity.

This is a question about Connes embedding~\cite{connes,hm11,hm15}. Haagerup and Musat showed
that every factorizable channel is a limit of finite matrix-tracial ones exactly when the
Connes embedding problem has a positive answer~\cite[Theorems~3.6 and~3.7]{hm15}; see
also~\cite[Corollary~3.4]{mr21} and~\cite[Corollary~3.11]{musat}. Since
$\mathrm{MIP}^*=\mathrm{RE}$~\cite{mip}, factorizable channels outside the closure exist, but the
argument does not produce one.

Nonhyperlinear groups have since been found through centralizer rigidity. Building on
OpenAI's nonsofic group~\cite{openai}, Kun and Thom showed that the amalgamated double $G*_HG$
of Kazhdan groups $H<G$, with $H$ infranormal but not normal, is nonsofic, by a
centralizer-normalization theorem~\cite[Theorems~A and~4.1]{kt}, and gave explicit such pairs
$\mathrm{EL}_r(\F_q[x_1,\ldots,x_d])<\mathrm{EL}_r(\F_q[x_1^{\pm1},\ldots,x_d^{\pm1}])\rtimes\SL_d(\Z)$
for every prime power $q$ and all $r,d\ge3$~\cite[Theorem~E]{kt}. Thom showed that a positive
answer to the centralizer problem of Alekseev and Thom~\cite[Open Problem~6.2(a)]{at} makes
these doubles nonhyperlinear, with an explicit witness commutator~\cite[Theorems~1.2 and~1.3,
Section~5.2]{thom}. Liu answered the problem~\cite[Theorem~1.2]{liu}; Alekseev, Liu and Thom
proved a stronger internality theorem, the form used here~\cite[Theorem~B]{alt}, and deduced
that the doubles are nonhyperlinear~\cite[Corollary~C]{alt}. Section~\ref{sec:group} states
the one consequence of these recent preprints~\cite{thom,alt} that we use.

A finite channel can enforce only finitely many relations, so it needs a finitely presented
group. These doubles are finitely generated and recursively presented, so Higman's embedding
theorem~\cite{higman} places them inside finitely presented groups, which are then
nonhyperlinear, but the embedding gives no usable presentation. We give an explicit one
instead: the amalgamated double of an explicitly presented Kazhdan group $\Ghat$ over a
finitely generated Kazhdan subgroup $\Hhat$, mapping into the double of the Kun--Thom pair with
$q=5$ and $r=d=3$. Property $(T)$ of $\Ghat$ and $\Hhat$ can be checked directly with the
three-subgroup criterion of Ershov and Jaikin-Zapirain~\cite{ejz}, and Thom's commutator
becomes a word of length six. Rational two-level blocks then turn the
presentation and the word into channel data, specified entry by entry. To our knowledge this
is the first explicit factorizable channel outside the closure of the finite matrix-tracial
factorizations.

It differs from earlier explicit examples in both directions. Nonfactorizable channels lie
outside the closure but have no tracial realization at all~\cite{hm11}; channels that need an
infinite ancilla for an exact factorization can still be approximated by finite
ones~\cite{mr,mori}. The channel here has an exact tracial factorization and no finite
matrix-tracial approximation. The appendices give the external input
(Appendix~\ref{sec:inputs}), the complete presentation and its concrete image
(Appendices~\ref{sec:presentation} and~\ref{sec:image}), the word arithmetic and deletion
certificates (Appendices~\ref{sec:arithmetic} and~\ref{sec:pruning}), the channels and the
global exclusion (Appendices~\ref{sec:channels} and~\ref{sec:global}), the quantitative bounds
and the memory argument (Appendices~\ref{sec:restricted}--\ref{sec:causal-full}), and the data
format (Appendix~\ref{sec:files}).

\paragraph{Notation.}
Write $M_d=M_d(\C)$ and $\tau_m=\Tr/m$. Let $T_d$ consist of the channels
\begin{equation}
 \Psi(X)=(\id\otimes\tau_m)\bigl[U(X\otimes I_m)U^*\bigr],
 \qquad U\in\U(dm),\quad m\ge1,
 \label{eq:finite-factor}
\end{equation}
and put $K_d=\overline{T_d}$. Finite convex mixtures have the same closure: rational mixture
weights are obtained by repeating matrix blocks, and arbitrary weights follow by
approximation~\cite[Proposition~3.8(ii)]{musat}. The class $\Fact_d$ of factorizable
channels, introduced by Anantharaman-Delaroche~\cite{ad} and characterized in the unitary
form used here by Haagerup and Musat~\cite[Theorem~2.2]{hm11}, allows a finite von Neumann
algebra with a faithful normal tracial state in place of $M_m$. Here ``finite'' is the
operator-algebraic notion; the algebra need not be finite-dimensional. In the notation of
Refs.~\cite{mr21,musat}, $T_d=\mathcal{FM}_{\mathrm{matr}}(d)$, its convex hull is
$\mathcal{FM}_{\mathrm{fin}}(d)$, and $\Fact_d=\mathcal{FM}(d)$. Every channel in $K_d$ is
factorizable (Lemma~\ref{lem:ultra-factor}), and by Ref.~\cite[Theorem~3.6]{hm15}, $K_d$ consists of
the factorizable channels with an ancilla that embeds trace-preservingly into the ultrapower
of the hyperfinite $\mathrm{II}_1$ factor, the qa-factorizable channels in the terminology of
Ref.~\cite{pr}.

We use the normalized Choi state and two channel distances,
\begin{align}
 J(\Phi)&=(\Phi\otimes\id)(\ket{\Omega_d}\bra{\Omega_d}),
 &\ket{\Omega_d}&=d^{-1/2}\sum_{k=0}^{d-1}\ket{kk},\\
 \dJ(\Phi,\Psi)&=\tfrac12\nnorm{J(\Phi)-J(\Psi)}_1,
 &\dd(\Phi,\Psi)&=\tfrac12\nnorm{\Phi-\Psi}_\diamond.
 \label{eq:distances}
\end{align}
Thus $\dJ\le\dd\le d\,\dJ$. Their distances to $K_d$ are denoted
$\Delta_J(\Phi)$ and $\Delta_\diamond(\Phi)$. All bath $2$-norms are normalized:
$\nnorm{x}_2^2=\tau(x^*x)$.

\begin{theorem}[Finite presentation and rational channels]\label{thm:main}
The finite rules in Section~\ref{sec:group} define a group $\Dhat$ with 43 generators and
4\,417 relators, and a named element $j\ne1$, such that every homomorphism from $\Dhat$ into a
tracial matrix ultraproduct sends $j$ to the identity. In particular, $\Dhat$ is
nonhyperlinear.

The channel construction in Section~\ref{sec:channel} gives channels $\Phi_1$, $\Phi_2$,
and $\Phi_{15}$ in dimensions $31\,981$, $31\,982$, and $32\,768$, respectively, with
\begin{equation}
 \Phi_\nu\in\Fact_{d_\nu}\setminus K_{d_\nu},
 \qquad \Delta_J(\Phi_\nu)>0.
 \label{eq:main-separation}
\end{equation}
Each has Choi rank $11\,230$ and Kraus entries in $\{0,1,3/5,4/5,-3/5\}$.
The channel $\Phi_{15}$ acts on exactly fifteen system qubits and has $64\,748$ nonzero
Kraus entries.
\end{theorem}

By Ref.~\cite[Theorem~3.6]{hm15}, no ancilla for these channels embeds trace-preservingly into
the ultrapower of the hyperfinite $\mathrm{II}_1$ factor; Section~\ref{sec:extraction} exhibits an
explicit ancilla, the separable $\mathrm{II}_1$ factor $L(\Gamma*F_2)$. Through its first part,
the whole theorem rests on the external theorem stated in Section~\ref{sec:group}; the
presentation, the
word calculations, the deleted equations and the rational matrices are verified by the
accompanying programs. We also keep the original fully referenced channel $\Phi_\full$ on
$M_{44086}$, whose extra reference positions serve the quantitative estimates of
Section~\ref{sec:quant}.

\section{A word that matrix approximations must lose}\label{sec:group}

A group is \emph{Kazhdan} if it has property $(T)$. For a subgroup $H<G$, its compression
semigroup is
\[
 P_H=\{t\in G:tHt^{-1}\le H\}.
\]
We use the following consequence of Theorem B of Ref.~\cite{alt} and Theorem 1.2 of
Ref.~\cite{thom}. Both are recent preprints, unrefereed at the time of writing, and this
statement is the only input taken from them:
\begin{quote}
If $H$ and $G$ are Kazhdan and $P_H$ generates $G$, then every homomorphism
$\pi:G\to\U(\M)$ into a tracial matrix ultraproduct satisfies
$\pi(g)(\pi(H)'\cap\M)\pi(g)^*=\pi(H)'\cap\M$ for every $g\in G$.
\end{quote}
No injectivity or trace-preservation hypothesis is imposed on $\pi$.

\subsection{The finite rules}
Let $C=(0,1,2,3,-1,-2,-3)$. Start with generators $f_r$, where
$f\in\{a,b,c\}$ and $r\in C$. Impose fifth powers, commutation within each family, and
\begin{equation}
 [[f_r,g_s],f_u]=[[f_r,g_s],g_u]=1
 \quad\bigl((f,g)=(a,b),(b,c),(c,a);\ r,s,u\in C\bigr),
 \label{eq:root-rel}
\end{equation}
with $[x,y]=xyx^{-1}y^{-1}$. Add six generators $t_{ij}$, $1\le i\ne j\le3$, with
\begin{equation}
 [t_{ij},t_{kl}]=1\ (i\ne l,\ j\ne k),
 \qquad [t_{ij},t_{jk}]=t_{ik}\ (i,j,k\text{ distinct}).
 \label{eq:steinberg}
\end{equation}
Trivial and duplicate commuting relations are omitted. For the cyclic successor and
predecessor $f^+$ and $f^-$ of $f$, set
\[
 B_f(r,s)=[[f_r,(f^+)_0],[(f^-)_s,f_0]].
\]
For $\epsilon=\pm1$, impose
\begin{equation}
 t_{ij}^{\epsilon}f_rt_{ij}^{-\epsilon}=
 \begin{cases}
 B_f(\epsilon i,j),&r=j,\\
 B_f(-\epsilon i,-j),&r=-j,\\
 f_r,&r=0\text{ or }|r|\ne j.
 \end{cases}
 \label{eq:conjugation}
\end{equation}
Finally put $\zeta_i=(t_{ik}t_{ki}^{-1}t_{ik})^2$, where $k=i+1$ cyclically, and impose
$f_{-i}=\zeta_i f_i\zeta_i^{-1}$. Call the resulting group $\Ghat$ and define its
\emph{actual subgroup}
\[
 \Hhat=\langle f_r:f\in\{a,b,c\},\ r=0,1,2,3\rangle<\Ghat.
\]

These finite rules supply the hypotheses of the external theorem. Any two root families generate
a finite class-two group. After removing common invariant vectors, their fixed spaces have
angle cosine at most $1/\sqrt5<1/2$. The three-subgroup criterion of Ref.~\cite{ejz}
therefore gives property $(T)$ for the root groups. The six action generators present
$\St_3(\Z)$, which is Kazhdan~\cite{ejz}. Both signs in
Eq.~\eqref{eq:conjugation} make the full root subgroup normal, and property $(T)$ passes to
extensions~\cite[Proposition~1.7.6]{bhv}, so $\Ghat$ is Kazhdan. The subgroup $\Hhat$ is a
quotient of the positive-root presentation and is Kazhdan as well. Each of the six $t_{ij}$,
taken with exponent $+1$, compresses $\Hhat$ into itself; the naming relations show that
they generate $\Ghat$ together with $\Hhat$.
Appendix~\ref{sec:presentation} gives the fixed-space calculation, the extension argument, and the
exact enumeration conventions.

\subsection{A concrete image and the surviving word}
Let
\[
 R_+=\F_5[x_1,x_2,x_3],\qquad
 R=\F_5[x_1^{\pm1},x_2^{\pm1},x_3^{\pm1}],\qquad
 G_0=\SL_3(R)\rtimes\SL_3(\Z).
\]
The second factor acts on exponent columns by monomial substitutions. Map $a_r,b_r,c_r$
to elementary matrices in positions $(1,2),(2,3),(3,1)$ with coefficient $z_r$, where
$z_0=1$ and $z_{\pm i}=x_i^{\pm1}$, and map $t_{ij}$ to the action matrix $I+E_{ij}$.
All defining relations hold. For $t=t_{12}$ and $h=a_2$, the image of $\Hhat$ has
polynomial entries, while the image of $t\Hhat t^{-1}$ has entries in
$\F_5[x_1,x_1x_2,x_3]$. The coefficient $x_2$ of $h$ is not in that subring. Thus
$h\notin t\Hhat t^{-1}$. This is the explicit compressor of Kun and Thom~\cite[proof of
Theorem~E]{kt}, also used by Thom~\cite[Section~5.2]{thom}.

Form $\Dhat=\Ghat*_{\Hhat}\Ghat$ by identifying the twelve positive-root generators.
Write $t_0,t_1$ for the two copies of $t$, and adjoin the redundant name
\begin{equation}
 j=[h,t_1t_0^{-1}],\label{eq:witness}
\end{equation}
the commutator from Thom's proof of his Theorem~1.3~\cite[Section~5.2]{thom}, written in
the presented group.
This gives 43 generators and 4\,417 distinct freely reduced relators, of total length
41\,953 and maximum length 19. The signed generator numbers of the witness word are
$(3,37,-22,-3,22,-37)$; its new generator number is 43.

The two concrete maps extend to
\[
 D_0=G_0*_{H_0}G_0,\qquad H_0=\SL_3(R_+)\times\{I\}.
\]
There the image of $j$ is
\begin{equation}
 (ht)_1\,(t^{-1}h^{-1}t)_0\,(t^{-1})_1.\label{eq:normal-word}
\end{equation}
The outside syllables have nonidentity monomial action. The middle one has coefficient
$-x_2/x_1$ in position $(1,2)$. All three lie outside $H_0$, so the reduced-word theorem
for an amalgam makes $j$ nontrivial; this is Thom's normal-form argument~\cite[Section~5.2]{thom}
carried out in $D_0$. Let $\Gamma$ be the image of $\Dhat$ in $D_0$.

Now, following Thom's proof~\cite[Section~5.2]{thom}, take any matrix-ultraproduct
representation $\pi$ of $\Dhat$ and put
$A=\pi(\Hhat)'\cap\M$. The element $c=t_0^{-1}t_1$ centralizes the shared subgroup:
both copies implement the same compression on it. Hence $\pi(c)\in A$. The external
normalization theorem gives
\[
 \pi(t_1t_0^{-1})=\pi(t_0ct_0^{-1})\in A.
\]
It commutes with $\pi(h)$, and $\pi(j)=I$. The same assertion holds for representations
of $\Gamma$ by composition with $\Dhat\to\Gamma$. No closure property of hyperlinear
groups under arbitrary quotients is needed.

\section{From the presentation to rational matrices}\label{sec:channel}

The computation uses $\Gamma$, not a presumed solution to the word problem of $\Dhat$.
Laurent-polynomial arithmetic decides multiplication in $G_0$ and membership in $H_0$.
Combining adjacent syllables and absorbing $H_0$-syllables decides the word problem in
$D_0$. To compare two words, reduce their quotient; reduced syllable lists need not be
unique.

Collect the identity, the signed generators, and the relator prefixes as labels, identifying
equal elements of $\Gamma$. A prefix step supplies an equation $xy=z$, represented by a
triple $(x,y,z)$. Include a generator-inverse equation for each positive generator, then
remove duplicate triples and equations with an identity input. The original table has
$|\Lab|=11\,230$ labels and $|\Tri_0|=16\,428$ triples. Label 0 is the identity; label
85 is the named element $j$.

\paragraph{Removing redundant equations.}
Retain all 43 positive-generator inverse equations. If $y^\vee$ denotes the label of the
inverse of a signed generator, then
\[
 xy=z,\quad yy^\vee=1\quad\Longrightarrow\quad zy^\vee=x.
\]
Whenever both multiplication equations occur, one is redundant. This removes 408 triples.
Another 30 triples merely reverse a retained inverse equation. The remaining table $\Tri$
has 15\,990 triples. Each deleted equation is supplied as a product of conjugates of
\emph{retained} equations; no deletion depends on another deletion. Every original relator
is therefore still enforced, and every label still occurs in the retained table.

\paragraph{The blocks.}
Choose a list $B$ of singleton reference labels containing 0. For each reference position
$p$ with label $g$, set $(A_g)_{pp}=1$. For the $n$th retained triple $(x,y,z)$, use positions
$p=|B|+2n,p+1$ and put
\begin{equation}
 (A_0)_{p,p}=\frac35,\qquad (A_y)_{p,p+1}=\frac45,\qquad
 (A_x)_{p+1,p}=\frac45,\qquad (A_z)_{p+1,p+1}=-\frac35.
 \label{eq:kraus-block}
\end{equation}
All other entries vanish. The channel is
\begin{equation}
 \Phi_B(X)=\sum_{g\in\Lab}A_gXA_g^*.\label{eq:channel}
\end{equation}
Taking $B=(0)$ gives $\Phi_1$; taking $B=(0,85)$ gives $\Phi_2$.
To obtain $\Phi_{15}$, append 786 identity singleton positions to $\Phi_2$, assigning all
of them to label 0. Equivalently, append 786 identity blocks to the factorizing unitary.
This prescription also fixes the action on off-diagonal matrix units involving the added
positions; it is not merely a specification of two diagonal channel corners.

\begin{table}[ht]
\centering
\small
\begin{tabular}{@{}lrrrr@{}}
\toprule
Channel & Dimension & Triples & Singletons & Nonzero entries\\
\midrule
Original $\Phi_\full$ & 44\,086 & 16\,428 & 11\,230 & 76\,942\\
One reference $\Phi_1$ & 31\,981 & 15\,990 & 1 & 63\,961\\
Two references $\Phi_2$ & 31\,982 & 15\,990 & 2 & 63\,962\\
Fifteen qubits $\Phi_{15}$ & 32\,768 & 15\,990 & 788 & 64\,748\\
\bottomrule
\end{tabular}
\caption{All four channels have Choi rank 11\,230 and common Kraus denominator five.
In the final row, 787 singleton positions have label 0 and one has label 85.}
\label{tab:instances}
\end{table}

\paragraph{General factorization and spectrum.}
Let $\lambda$ be the regular representation of $\Gamma$, with its canonical trace. The
block associated with $(x,y,z)$ is
\begin{equation}
 \begin{pmatrix}(3/5)I&(4/5)\lambda(y)\\(4/5)\lambda(x)&-(3/5)\lambda(z)\end{pmatrix}
 =\operatorname{diag}(I,\lambda(x))
 \left(\begin{pmatrix}3/5&4/5\\4/5&-3/5\end{pmatrix}\otimes I\right)
 \operatorname{diag}(I,\lambda(y)).\label{eq:unitary-block}
\end{equation}
It is unitary because $xy=z$. Its direct sum with the singleton blocks is
$U=\sum_gA_g\otimes\lambda(g)$. Distinct group elements are orthogonal in the canonical
trace, so $U$ realizes Eq.~\eqref{eq:channel} tracially; this is the factorization
criterion of Ref.~\cite[Theorem~2.2]{hm11}, up to their convention of conjugating by $u^*$. In particular, the channel is unital and trace
preserving.

The Kraus supports are nonempty and pairwise disjoint. Thus the Choi rank is $|\Lab|$, and
its nonzero eigenvalues are
\begin{equation}
 \frac{b_g}{25d},\qquad
 b_g=\sum_{(p,q,n)\in\mathrm{kraus}[g]}n^2.
 \label{eq:choi-spectrum}
\end{equation}
For $\Phi_{15}$ the common spectral denominator is 819\,200. The package includes the
complete weight histograms. These are exact spectral data, not estimates of the global
separation.

\section{Why one identity reference is enough}\label{sec:extraction}

The original construction gave every label a singleton. That makes extraction transparent,
but it is unnecessary for an exact realization. Once the identity label is unitary, the row
and column norms of a two-level block force the other labels to be unitary too.

\begin{lemma}[Extraction from Choi support]\label{lem:extract}
Let $\Phi_B$ be a channel constructed above, and let $\Pi=\supp J(\Phi_B)$.
If a unitary $V\in M_d\otimes N$ in a finite von Neumann algebra with faithful tracial
state implements a channel $\Psi$ satisfying $\supp J(\Psi)\le\Pi$, then
\[
 V=\sum_g A_g\otimes v_g
\]
for unitaries $v_g\in N$. The normalized labels $w_g=v_0^*v_g$ satisfy $w_0=I$ and
every retained multiplication equation. If $\Psi=\Phi_B$, then also
$\tau(v_gv_h^*)=\delta_{g,h}$.
\end{lemma}
\begin{proof}
The range of $\Pi$ is the span of the vectorized Kraus matrices. Every vector orthogonal
to that span has zero quadratic form in $J(\Psi)$; the corresponding linear combination
of the bath blocks of $V$ has zero $2$-norm and vanishes. Therefore $V=\sum_gA_g\otimes v_g$.

The identity singleton makes $v_0$ unitary. Write $a=3/5$ and $b=4/5$. Each two-level
block of $V$ has the form
\[
 \begin{pmatrix}av_0&bv_y\\bv_x&-av_z\end{pmatrix}.
\]
Its first row and first column imply $v_yv_y^*=I$ and $v_x^*v_x=I$.
In a finite tracial algebra an isometry or coisometry is unitary. The second column then
gives $v_z^*v_z=I$. Every label occurs in a block or a singleton, so all $v_g$ are unitary.
Row orthogonality gives
\[
 v_0v_x^*=v_yv_z^*,\qquad v_z=v_xv_0^*v_y,
\]
and hence $w_xw_y=w_z$. Finally, when the Choi states are equal, comparison of coefficients
in the orthogonal basis $(\vket{A_g})_g$ gives the stated trace identities, as in
Ref.~\cite[Theorem~2.2]{hm11}.
\end{proof}

\begin{proposition}[Presentation-to-channel principle]\label{prop:compiler}
Suppose a finite presentation has a named word $j$ killed by every matrix-ultraproduct
representation, and a concrete quotient in which $j\ne1$. Take as labels finitely many
pairwise distinct elements of that quotient, including the identity, the images of the
generators and of their inverses, and the image of $j$, and as triples multiplication
equations $xy=z$ that hold there. The resulting channel lies in
$\Fact_d\setminus K_d$ whenever the retained equations imply every presentation relator and
generator-inverse equation and identify the label of $j$ with the word $j$, every label is
covered by a block or singleton, and an identity singleton is present.
\end{proposition}
\begin{proof}
General factorization is Eq.~\eqref{eq:unitary-block}; it needs the labels to be distinct
and the triples to be valid in the quotient. If finite matrix-tracial channels converged to
the target, their bath-block ultraproduct would give an exact unitary factorization of the
target in a tracial matrix ultraproduct. Compare the implication (2)$\Rightarrow$(1) of
Ref.~\cite[Theorem~3.6]{hm15}, which is proved there through condition (3), with a
factorization through the ultrapower of the hyperfinite factor; here the matrix ultraproduct
is taken directly. The system dimension is fixed, so there are only finitely many blocks and
Choi entries to pass to the limit.

Lemma~\ref{lem:extract} gives unitary labels satisfying the presentation. The representation
therefore sends $j$ to $I$. But the identity and $j$ are distinct labels, so
\[
 1=\tau(w_j)=\tau(v_0^*v_j)=0,
\]
a contradiction. The approximating unitaries were never assumed to have the prescribed
block form; that form is forced only in the exact limit.
\end{proof}

Encoding operator relations in second moments has precedents. Musat and R\o rdam force
linear relations among unitaries through their trace correlations, following Regev,
Slofstra, and Vidick~\cite[Proposition~3.2]{mr}, and turn them into a Schur-multiplier
channel that needs a type $\mathrm{II}_1$ ancilla. Mori uses correlation matrices of unitaries
to force a product relation among symmetries that has no finite-dimensional
solution~\cite{mori}. Passing from a
presentation to an object in which a killed word must survive is Slofstra's method for
nonlocal games~\cite{slofstra}; see also Ref.~\cite{pr} for games and factorizable maps.
Proposition~\ref{prop:compiler} applies these ideas to group relations in a channel.

The deletion certificates verify the hypotheses of Proposition~\ref{prop:compiler} for all
three reduced channels. Closedness of $K_d$ gives a positive distance. This completes the
proof of Theorem~\ref{thm:main}.

Replacing the bath $L(\Gamma)$ by $L(\Gamma*F_2)$, where $F_2$ is the free group on two
generators, preserves the trace table and gives a separable $\mathrm{II}_1$ factor
(Appendix~\ref{sec:global}). Since $\Phi_\nu\notin K_{d_\nu}$, the implication (1)$\Rightarrow$(2) of
Ref.~\cite[Theorem~3.6]{hm15} shows that this factor is not Connes-embeddable; compare the
existential version in Ref.~\cite[Corollary~3.11]{musat}. At the level of groups this is the standard correspondence between hyperlinear groups and
embeddable group von Neumann algebras~\cite[Proposition~7.1]{ozawa}. Liu's manuscript draws
the corresponding non-embedding for the group von Neumann algebra of a wreath product from its
Theorem~1.1~\cite{liu}.

\section{An elementary separator on the exact support face}\label{sec:support}

The global obstruction concerns limits of matrix representations. A simpler argument
already explains what happens to \emph{exact} finite representations, without property $(T)$
or ultraproduct rigidity: the finite-dimensional form of Thom's observation that a compressor
always normalizes a finite image of the subgroup~\cite[Remark~4.4]{thom}, combined with his
argument for the witness~\cite[Section~5.2]{thom}.

\begin{lemma}[Finite-dimensional commutants]\label{lem:finite}
If $H<G$, $tHt^{-1}\le H$, and $D=G*_HG$, then every homomorphism
$\pi:D\to\U(m)$ kills $[h,t_1t_0^{-1}]$ for every $h\in H$.
\end{lemma}
\begin{proof}
Put $A=\pi(H)'\cap M_m$. Compression of $H$ gives
$\pi(t_0)A\pi(t_0)^*\supseteq A$. Conjugation preserves vector-space dimension, so this
inclusion is equality. As before, $t_0^{-1}t_1$ centralizes $H$; thus
$\pi(t_1t_0^{-1})\in A$ and it commutes with $\pi(h)$.
\end{proof}

For a target channel with support projection $\Pi$, define
\[
 T_d^\Pi=\{\Psi\in T_d:\supp J(\Psi)\le\Pi\}.
\]
This is a restricted comparison class. Its closure is not being identified with
$K_d\cap\{\supp J\le\Pi\}$.

\begin{theorem}[Support-restricted separators]\label{thm:support}
For any channel in Table~\ref{tab:instances}, put $s_g=\nnorm{A_g}_{\mathrm{HS}}^2$.
Then
\begin{equation}
 \inf_{\Psi\in\overline{T_d^\Pi}}\dJ(\Phi,\Psi)
 \ge \frac{\sqrt{s_0s_{85}}}{d}.
 \label{eq:restricted-choi}
\end{equation}
For $\Phi_\full$, $\Phi_2$, and $\Phi_{15}$, which retain separate identity and witness
singletons, one also has
\begin{equation}
 \inf_{\Psi\in\overline{T_d^\Pi}}\dd(\Phi,\Psi)\ge\frac12.
 \label{eq:restricted-diamond}
\end{equation}
The last bound is witnessed by one unentangled input and one two-outcome measurement.
\end{theorem}
\begin{proof}
Let $\Psi\in T_d^\Pi$. Lemma~\ref{lem:extract} gives unitary labels satisfying the retained
equations. The deletion certificates (Proposition~\ref{prop:pruning}) and the prefix lemma
(Lemma~\ref{lem:prefix}) turn these into
a representation of $\Dhat$, and Lemma~\ref{lem:finite} then gives $v_{85}=v_0$.
The normalized vectors $\ket{a_g}=\vket{A_g}/\sqrt{s_g}$ are orthogonal. The target Choi
state has zero cross entry between $\ket{a_0}$ and $\ket{a_{85}}$; the comparison has
cross entry $\sqrt{s_0s_{85}}/d$, while their two diagonal entries agree.
The projection onto $(\ket{a_0}+\ket{a_{85}})/\sqrt2$ proves
Eq.~\eqref{eq:restricted-choi}.

Let $p_0,p_j$ be the two singleton system positions and input
$\ket{+}=(\ket{p_0}+\ket{p_j})/\sqrt2$. The target output is the equal incoherent mixture
of those positions. Since $v_j=v_0$, every supported comparison leaves $\ket{+}\bra{+}$
unchanged. Their output half trace distance is $1/2$; measuring $\ket{+}\bra{+}$ attains
this difference. Continuity extends both bounds to the closure of the restricted class.
\end{proof}

Rational mixtures of supported finite factorizations are again supported finite
factorizations, obtained by repeating blocks, so both bounds hold against rational mixtures
and, by approximation, against all finite mixtures, which lie in the closure.

For $\Phi_{15}$, the integer weights are $b_0=163\,972$ and $b_{85}=41$.
Equation~\eqref{eq:restricted-choi} is therefore strictly larger than $1/316$, as checked by
squaring and comparing integers. The corresponding rational lower bounds for $\Phi_1$
and $\Phi_2$ are $1/527$ and $1/329$.

These numerical separators do not evaluate $\Delta_J$ or $\Delta_\diamond$. A nearby
finite-tracial channel can have small nonzero Choi weight outside $\Pi$. Projecting that
weight away is not known to preserve a unitary finite-tracial realization. The support
hypothesis cannot be removed from Theorem~\ref{thm:support} by continuity alone.

\section{Quantitative extraction and the remaining tolerance}\label{sec:quant}

The universal-kernel result has a uniform finite-error consequence. There is a rational
$e_0\in(0,1]$ such that, in every matrix dimension and for every unitary assignment to the
43 positive generators,
\begin{equation}
 \max_{r\in\Rel}\nnorm{\varphi(r)-I}_2\le e_0
 \quad\Longrightarrow\quad \nnorm{\varphi(j)-I}_2\le\frac12.
 \label{eq:tolerance}
\end{equation}
Otherwise assignments with vanishing relator defects and witness displacement exceeding
$1/2$ would yield a contradicting ultraproduct representation. This is the standard
compactness argument; Liu's manuscript draws the same kind of uniform consequence for its
group~\cite[Theorem~1.1]{liu}. It proves existence, not a numerical value of $e_0$.

\begin{theorem}[Global bounds from a relator tolerance]\label{thm:quant}
For the original fully referenced channel $\Phi_\full$, every tolerance satisfying
Eq.~\eqref{eq:tolerance} gives
\begin{align}
 \Delta_J(\Phi_\full)&\ge\frac{e_0^2}{65\,188\,028\,416},
 \label{eq:q-choi}\\
 \Delta_\diamond(\Phi_\full)&\ge\frac{e_0^2}{1\,478\,656}.
 \label{eq:q-diamond}
\end{align}
More generally, for a fully referenced channel built by the same block rule, the
denominators are $4096\ell_*^2d$ and $4096\ell_*^2$ when the maximum relator length is
$\ell_*\ge2$ and every relator path and generator inverse is encoded by the triangle table.
\end{theorem}

The bounds come from the following extraction; Appendix~\ref{sec:quantitative} gives each estimate. Let $U$
implement an arbitrary finite-tracial comparison. A singleton bath block $B_g$ is a
contraction. Small probability of leaving its system position gives a unitary polar
completion $v_g$ with $\nnorm{B_g-v_g}_2\le k$, where $k=\sqrt{d\delta}$ for Choi error
$\delta$, or $k=\sqrt\eta$ for half-diamond error $\eta$. A zero-target Gram constraint
then places each triangle entry within $3k$ of its prescribed scalar multiple of $v_g$.
In the half-diamond case this uses a probe supported on two input positions, rather than
the uniform $d$-dimensional Choi probe.

Unitarity controls the row product from outside each two-level block. Row orthogonality
therefore gives triangle defects below $32k$, and following the original relator paths gives
\[
 \max_{r\in\Rel}\nnorm{\varphi(r)-I}_2\le64\ell_*k.
\]
The identity and witness references also give
$|\tau(v_0^*v_j)|\le2k^2+2k$. At $k\le e_0/(64\ell_*)$, this contradicts
Eq.~\eqref{eq:tolerance}. The argument covers arbitrary finite baths, with no support
restriction.

The smaller channels do not automatically inherit these constants: they no longer have a
separate singleton for every label. We retain $\Phi_\full$ precisely to keep this
quantitative route available.

A Kazhdan angle constant cannot simply be substituted for $e_0$. For example, assign one
generator $\operatorname{diag}(1,z_m)$ and all others $I_2$, where
\[
 z_m=\frac{m^2-1+2mi}{m^2+1}.
\]
The relator defects tend to zero, but the lazy conjugation gap off the coordinate commutant
is $1/[43(m^2+1)]$. The witness also approaches the identity, so this is not a
counterexample to the group theorem. It shows why quantitative lifting or normalization,
not the exact group's spectral gap alone, is needed.

\section{Purity and memory in causal repeated use}\label{sec:causal}

The separation also has a resource interpretation. Use the causal model of
Ref.~\cite[Definition~2.3]{causal}: each output is released before the next input arrives; all retained
registers count toward a memory of $Q$ qubits; imported states are fixed and independent;
and discarded registers never return. Write $P=Q-S(\beta_0)$ for initial purity,
$C$ for the sum of imported qubits minus their entropy, and $J$ for the number of
replacement slots. An export and its replacement count once. Error $\epsilon$ is the
largest half trace distance of the complete experiment, over adaptive observers with
quantum references.

For any channel in Table~\ref{tab:instances}, the purity converse~\cite[Theorem~3.2(b)]{causal} gives
\begin{equation}
 P+C\ge\frac{n(\Delta_J(\Phi)-\epsilon)_+^2}{2\ln2}.
 \label{eq:purity}
\end{equation}
Appendix~\ref{sec:causal-full} reproduces the local tracialization and entropy ledger needed for this
inequality. It imposes no restriction on $Q$ or $J$.

The fully referenced channel has a pure input-reference probe whose output is flat of
rank $r=11\,230$; it also has maximal complementary entropy $\log_2r$. Proposition 8.7
of Ref.~\cite{configuration} sharpens the earlier flat-probe converse~\cite[Theorem~5.8]{causal} to $P+C<3Q+6$ at
error at most $1/16$ and exchange $J\le n\log_2r/2$. Combining the two inequalities yields
\begin{equation}
 Q\ge\frac{n(\Delta_J(\Phi_\full)-\epsilon)_+^2}{6\ln2}-2.
 \label{eq:memory-general}
\end{equation}
In particular,
\begin{equation}
 Q\ge\frac{n\Delta_J(\Phi_\full)^2}{24\ln2}-2,
 \qquad \epsilon\le\min\{1/16,\Delta_J(\Phi_\full)/2\}.
 \label{eq:memory}
\end{equation}
The exchange restriction is not claimed optimal outside $K_d$.

The two-reference channels $\Phi_2$ and $\Phi_{15}$ retain a flat rank-two probe:
the input $\ket{+}$ of Theorem~\ref{thm:support}. Thus
Eq.~\eqref{eq:memory-general} holds for each of them with its own $\Delta_J$ under the
more restrictive budget $J\le n/2$. We do not transfer the original rank-$11\,230$ probe
to a channel from which its reference positions have been removed.

\section{Data and open questions}

The package contains the literal presentation, the original label table, all reduced sparse
Kraus matrices, and a free-group certificate for every deleted multiplication equation. The
reduced instances are reconstructed independently, by deleting and reindexing slots of the
original channel. The checks cover all labels, all retained equations, both Kraus
normalization identities, the Choi weights and the integer arithmetic in the restricted
bounds, and deliberate corruptions are rejected. They do not formalize the external rigidity
theorems or the analytical estimates; Appendix~\ref{sec:files} records which statements are
finite checks and which rest on external theorems.

Two questions remain. The first is a number: an evaluated relator tolerance in
Eq.~\eqref{eq:tolerance} would bound $\Delta_J(\Phi_\full)$ explicitly, and a quantitative
support-leakage theorem would do the same for the reduced channels, whereas $1/316$ and $1/2$
hold only against supported comparisons. The second is size: fifteen system qubits are not
claimed to be minimal, and no efficient gate implementation is given. A finite pure environment
still realizes the channel exactly, while by Eq.~\eqref{eq:purity} with $n=1$ any finite
environment that realizes it to within $\epsilon$ must start with purity at least
$(\Delta_J-\epsilon)_+^2/(2\ln2)$ bits; the maximally mixed environment is the extreme case.

\paragraph{AI assistance.}
Generative AI was used in mathematical exploration, code development, checking and writing.
The authors are responsible for the content.

\clearpage
\appendix
\counterwithin{theorem}{section}
\counterwithin{equation}{section}
\counterwithin{table}{section}
% Appendices: the full finite presentation, the exact algorithms and all proofs.
\section{Conventions and external inputs}\label{sec:inputs}

Products compose from right to left and $[x,y]=xyx^{-1}y^{-1}$. The normalized trace on
$M_m(\C)$ is $\tau_m=\Tr/m$; $\nnorm{x}_2^2=\tau_m(x^*x)$, while
$\nnorm{x}_{\mathrm{HS}}^2=\Tr(x^*x)$ denotes the unnormalized Hilbert--Schmidt norm.
For a fixed nonprincipal ultrafilter $\omega$ on $\N$, a tracial matrix ultraproduct is
\begin{equation}
 \M=\prod_\omega(M_{m_n},\tau_{m_n})
 =\frac{\{(x_n):\sup_n\nnorm{x_n}_\infty<\infty\}}
 {\{(x_n):\lim_\omega\nnorm{x_n}_{2,m_n}=0\}}.
 \label{eq:ultra}
\end{equation}
It has faithful trace $\tau((x_n)_\omega)=\lim_\omega\tau_{m_n}(x_n)$.
A group is hyperlinear if it admits an injective homomorphism into the unitary group of
such an ultraproduct. Our universal-kernel statement quantifies over \emph{all}
homomorphisms, not just injective or trace-preserving ones.

The property-$(T)$ inputs are the three-subgroup criterion and the Steinberg-group theorem
of Ref.~\cite{ejz}. For the matrix-ultraproduct step we use the following external result.

\begin{theorem}[{External normalization input: Ref.~\cite[Theorem~1.2]{thom} with Ref.~\cite[Theorem~B]{alt}}]\label{thm:external}
Let $H<G$ be Kazhdan groups and suppose that
$P_H=\{t\in G:tHt^{-1}\le H\}$ generates $G$. For every homomorphism
$\pi:G\to\U(\M)$,
\begin{equation}
 \pi(g)(\pi(H)'\cap\M)\pi(g)^*=\pi(H)'\cap\M
 \qquad(g\in G).
 \label{eq:normalization}
\end{equation}
\end{theorem}

Theorem B of Ref.~\cite{alt} supplies internality of the relevant relative commutants;
Theorem 1.2 of Ref.~\cite{thom} then gives Eq.~\eqref{eq:normalization}. Property $(T)$
supplies the conjugation spectral gap needed for internality. Neither finite testing nor
same-dimension repair of approximate representations is substituted for them.

The centralizer problem answered in Ref.~\cite{alt} is Open Problem 6.2(a) of Ref.~\cite{at};
the proofs in Ref.~\cite{alt} are based on Liu's manuscript, whose Theorem~1.2 answers it~\cite{liu}.
Liu's manuscript is not on arXiv and is cited from an archived copy for provenance; the input
used here is Theorem~B of Ref.~\cite{alt}. Neither Theorem~B of Ref.~\cite{alt} nor
Theorem~1.2 of Ref.~\cite{thom} assumes that the dimensions $m_n$ grow. If they are bounded along $\omega$, the
ultraproduct is a matrix algebra, and the conclusion used below, $\pi(j)=I$, also follows
from Lemma~\ref{lem:finite-rigid}.

\section{The literal finite presentation}\label{sec:presentation}

\subsection{Three root families}
For an ordered finite set $C$ of size $m$, let $\Delta(C)$ have generators
$f_r$ with $f\in\{a,b,c\}$ and $r\in C$, subject to
\begin{align}
 f_r^5&=1,\label{eq:fifth}\\
 [f_r,f_s]&=1\quad(r<s\text{ in the chosen order}),\label{eq:within}\\
 [[f_r,g_s],f_u]&=[[f_r,g_s],g_u]=1
 \quad\bigl((f,g)=(a,b),(b,c),(c,a);\ r,s,u\in C\bigr).
 \label{eq:class-two}
\end{align}
The order in Eq.~\eqref{eq:within} fixes enumeration only.

\begin{lemma}\label{lem:kazhdan-roots}
Every $\Delta(C)$ has property $(T)$.
\end{lemma}
\begin{proof}
A single family generates an elementary abelian group of exponent five. For consecutive
families $X=\langle f_r\rangle$ and $Y=\langle g_s\rangle$, the mixed commutators are
central in $\langle X,Y\rangle$ by Eq.~\eqref{eq:class-two}, and
$[f_r,g_s]^5=[f_r^5,g_s]=1$. Collection expresses every element using at most
$2m+m^2$ generators or central commutators, with exponents from $0$ to $4$. Thus the
universal two-family group has order at most $5^{2m+m^2}$. Additional relations elsewhere
can only produce a quotient of this finite group.

Remove the invariant part of a unitary representation of the two-family group and decompose
it into irreducibles. In an irreducible summand, central commutators act by scalars, defining
a bicharacter
\[
 \chi:X\times Y\longrightarrow\{z:z^5=1\}.
\]
If $\chi$ is trivial, the summand is a nontrivial character of an abelian quotient, and at
least one of its $X$- or $Y$-fixed spaces is zero. Otherwise
\[
 Y_0=\{y\in Y:\chi(x,y)=1\text{ for every }x\in X\}
\]
has index at least five. If $v$ is $X$-fixed and $y\notin Y_0$, conjugation by a suitable
$x$ multiplies $\langle v,yv\rangle$ by a nontrivial fifth root of unity while leaving it
unchanged. Hence this coefficient vanishes. Averaging over $Y$ gives
\[
 \nnorm{P_Yv}^2=\langle v,P_Yv\rangle
 \le\frac{|Y_0|}{|Y|}\nnorm v^2\le\frac15\nnorm v^2.
\]
The same bound holds in a general representation without nonzero
$\langle X,Y\rangle$-invariant vectors. The projections $P_X$ and $P_Y$ commute with the
projections onto the irreducible summands, so an $X$-fixed vector is an orthogonal sum
$v=\sum_iv_i$ of $X$-fixed vectors in the summands, and
$\nnorm{P_Yv}^2=\sum_i\nnorm{P_Yv_i}^2\le\frac15\nnorm v^2$. This is the supremum over
summands used in the proof of Ref.~\cite[Corollary~4.7]{ejz}, and removing the common
invariant part restricts to the representations considered there.
The cosine between the two fixed spaces, after removing their common invariant part, is
therefore at most $1/\sqrt5$. Lifting representations preserves the bound in every
quotient. Relative property $(T)$ with respect to each finite family is automatic. The
three-subgroup criterion~\cite{ejz}, with $1/\sqrt5<1/2$, proves the claim.
\end{proof}

\subsection{Action generators and conjugation rules}
Let $T$ have six generators $t_{ij}$ for $1\le i\ne j\le3$, with relations
\begin{align}
 [t_{ij},t_{kl}]&=1\quad(i\ne l,\ j\ne k),\label{eq:t-comm}\\
 [t_{ij},t_{jk}]&=t_{ik}\quad(i,j,k\text{ distinct}).\label{eq:t-prod}
\end{align}
After omitting same-generator commutators and duplicate unordered commuting pairs, there
are six commuting relations and six product relations.

\begin{lemma}\label{lem:steinberg}
These relations present $T\cong\St_3(\Z)$, and $T$ is Kazhdan.
\end{lemma}
\begin{proof}
Put $E_{ij}(n)=t_{ij}^n$. Additivity is immediate. Equation~\eqref{eq:t-comm} gives the
commuting-root relations at all integer powers. For distinct $i,j,k$, the element $t_{ik}$
commutes with both $t_{ij}$ and $t_{jk}$, so
\[
 [t_{ij}^u,t_{jk}^v]=t_{ik}^{uv}\qquad(u,v\in\Z).
\]
These are the remaining Steinberg relations. The maps between $T$ and $\St_3(\Z)$
are inverse on generators. Property $(T)$ follows from Ref.~\cite{ejz}. No identification
of $\St_3(\Z)$ with $\SL_3(\Z)$ is used.
\end{proof}

Fix $C=(0,1,2,3,-1,-2,-3)$ and start with $\Delta(C)$ and $T$. For cyclic successor and
predecessor families $f^+,f^-$, define the word
\begin{equation}
 B_f(r,s)=[[f_r,(f^+)_0],[(f^-)_s,f_0]].\label{eq:bword}
\end{equation}
For every $i\ne j$, $\epsilon\in\{1,-1\}$, $f\in\{a,b,c\}$, and $r\in C$, impose
\begin{equation}
 t_{ij}^{\epsilon}f_rt_{ij}^{-\epsilon}=
 \begin{cases}
 B_f(\epsilon i,j),&r=j,\\
 B_f(-\epsilon i,-j),&r=-j,\\
 f_r,&r=0\text{ or }|r|\ne j.
 \end{cases}\label{eq:conj-full}
\end{equation}
The signed indices in $B_f$ are coefficient labels, not generator powers. Let $k=i+1$
cyclically modulo three, set $\zeta_i=(t_{ik}t_{ki}^{-1}t_{ik})^2$, and add
\begin{equation}
 f_{-i}=\zeta_i f_i\zeta_i^{-1}
 \qquad(f\in\{a,b,c\},\ i=1,2,3).
 \label{eq:negative-names}
\end{equation}
Call the resulting group $\Ghat$ and define
\begin{equation}
 \Hhat=\langle a_r,b_r,c_r:r=0,1,2,3\rangle<\Ghat.
 \label{eq:hhat}
\end{equation}

\begin{proposition}\label{prop:kazhdan-pair}
The groups $\Hhat$ and $\Ghat$ are Kazhdan. The six $t_{ij}$ compress $\Hhat$ into
itself and generate $\Ghat$ together with $\Hhat$.
\end{proposition}
\begin{proof}
The subgroup $\Hhat$ is a quotient of $\Delta(0,1,2,3)$, so it is Kazhdan by
Lemma~\ref{lem:kazhdan-roots}. The subgroup $N$ generated by all 21 root symbols is a
Kazhdan quotient of $\Delta(C)$. Both signs in Eq.~\eqref{eq:conj-full} make $N$
normal in $\Ghat$. The quotient $\Ghat/N$ is a quotient of $T$, so it has property
$(T)$ by Lemma~\ref{lem:steinberg}. Property $(T)$ passes to
extensions~\cite[Proposition~1.7.6]{bhv}, so $\Ghat$ has it.

For clarity, here is that argument for ordinary unitary representations. Property
$(T)$ of the normal subgroup forces almost invariant unit vectors to approach its invariant
subspace. This subspace is invariant under $\Ghat$, and the quotient has property $(T)$,
so it contains a nonzero $\Ghat$-invariant vector.

For $\epsilon=1$ and $r\ge0$, the right side of Eq.~\eqref{eq:conj-full} uses only
positive labels. Thus $t_{ij}\Hhat t_{ij}^{-1}\le\Hhat$. Equation~\eqref{eq:negative-names}
recovers every negative root symbol from the positive ones and the action generators.
These elements generate $\Ghat$. Since $\Hhat\subseteq P_{\Hhat}$, the whole compression
semigroup generates $\Ghat$ as well.
\end{proof}

\section{The computable image and the universal-kernel word}\label{sec:image}

Let
\begin{equation}
 R_+=\F_5[x_1,x_2,x_3],\qquad R=\F_5[x_1^{\pm1},x_2^{\pm1},x_3^{\pm1}].
\end{equation}
For $A\in\SL_3(\Z)$ define $\sigma_A(x^v)=x^{Av}$, with exponents written as columns.
The semidirect product law is
\begin{equation}
 G_0=\SL_3(R)\rtimes_\sigma\SL_3(\Z),\qquad
 (M,A)(N,B)=(M\sigma_A(N),AB).
 \label{eq:semidirect}
\end{equation}
With $e_{ij}(z)=I+zE_{ij}$, set
\begin{align}
 a_r&\longmapsto(e_{12}(z_r),I),&
 b_r&\longmapsto(e_{23}(z_r),I),&
 c_r&\longmapsto(e_{31}(z_r),I),\label{eq:root-images}\\
 t_{ij}&\longmapsto(I,I+E_{ij}),&&&\label{eq:t-images}
\end{align}
where $z_0=1$ and $z_{\pm i}=x_i^{\pm1}$.

\begin{lemma}\label{lem:image}
Equations~\eqref{eq:root-images}--\eqref{eq:t-images} define a homomorphism
$\Ghat\to G_0$. For $t=t_{12}$ and $h=a_2$, one has
$h\in\Hhat\setminus t\Hhat t^{-1}$.
\end{lemma}
\begin{proof}
Fifth powers vanish in characteristic five, and consecutive elementary-root families have
central mixed commutators. The integer action matrices satisfy the Steinberg relations.
For example, the image of $B_a(r,s)$ is
\[
 [e_{13}(z_r),e_{32}(z_s)]=e_{12}(z_rz_s);
\]
the other families follow by cycling matrix indices. The action $I+\epsilon E_{ij}$
sends $x_j$ to $x_i^\epsilon x_j$ and $x_j^{-1}$ to $x_i^{-\epsilon}x_j^{-1}$,
which proves Eq.~\eqref{eq:conj-full}. The image of $\zeta_i$ is minus identity on
coordinates $i,k$ and identity on the remaining coordinate; it inverts $x_i$, proving
Eq.~\eqref{eq:negative-names}.

The image of $\Hhat$ has polynomial entries. Its conjugate by $t_{12}$ has entries in
$\F_5[x_1,x_1x_2,x_3]$. Every monomial in this subring has first exponent at least its
second. The coefficient $x_2$ in the image of $h=a_2$ has exponent $(0,1,0)$ and is not
in that subring. This is the compressor in the proof of Theorem~E of Ref.~\cite{kt}, also
used in Ref.~\cite[Section~5.2]{thom}. Faithfulness of the homomorphism is not required.
\end{proof}

\subsection{The double and its enumeration}
Form
\[
 \Dhat=\Ghat*_{\Hhat}\Ghat
\]
by taking two copies of the presentation and identifying the twelve positive-root generators.
The subgroup $\Hhat$ need not have a separate finite presentation: the amalgamated-product
universal property only requires its finite generating set for these identifications.
Adjoin a named generator $j$ with relator
\begin{equation}
 j^{-1}[h,t_1t_0^{-1}]=1.
 \label{eq:j-name}
\end{equation}
This is a Tietze transformation and does not change the group. The commutator
$[h,t_1t_0^{-1}]$ and the normal-form argument below are those of Thom's proof of his
Theorem~1.3~\cite[Section~5.2]{thom}; here they are carried out in the presented group and in
its concrete image.

The concrete maps agree on $\Hhat$, giving a map into
\begin{equation}
 D_0=G_0*_{H_0}G_0,\qquad H_0=\SL_3(R_+)\times\{I\}.
 \label{eq:concrete-double}
\end{equation}
The image of $j$ has the reduced form
\begin{equation}
 (ht)_1\,(t^{-1}h^{-1}t)_0\,(t^{-1})_1.
 \label{eq:j-normal}
\end{equation}
The first and last syllables have nonidentity action matrices. The middle one has matrix
$I-(x_2/x_1)E_{12}$ and lies outside $H_0$. Thus all syllables are outside the amalgamating
subgroup and their sides alternate. The word is nontrivial. In particular $j\ne1$ in
$\Dhat$ and in $\Gamma=\operatorname{im}(\Dhat\to D_0)$.

\begin{table}[ht]
\centering\small
\begin{tabular}{@{}rl@{}}
\toprule
Indices & Generators in enumeration order\\
\midrule
1--12 & $a_0,a_1,a_2,a_3,b_0,b_1,b_2,b_3,c_0,c_1,c_2,c_3$\\
13--21 & $a_{-1},a_{-2},a_{-3},b_{-1},b_{-2},b_{-3},c_{-1},c_{-2},c_{-3}$, copy 0\\
22--27 & $t_{12},t_{13},t_{21},t_{23},t_{31},t_{32}$, copy 0\\
28--36 & The nine negative-root symbols in the same order, copy 1\\
37--42 & The six action generators in the same order, copy 1\\
43 & $j$\\
\bottomrule
\end{tabular}
\caption{A negative integer denotes the inverse of a generator, not a negative coefficient
label. For example, generator 14 is $a_{-2}$ in copy 0, whereas $-3$ is the inverse of $a_2$.}
\end{table}

The witness word is $(3,37,-22,-3,22,-37)$ and its defining relator is
$(-43,3,37,-22,-3,22,-37)$. Before deduplication, each copy contributes
\[
 (21+63+2058)+12+252+9=2415
\]
relators. The copies share $12+18+384=414$ positive-root relations. Including the new
name gives $2\cdot2415-414+1=4417$ distinct freely reduced relators.

\begin{table}[ht]
\centering\small
\begin{tabular}{@{}lr@{}}
\toprule
Family & Number retained\\
\midrule
Fifth powers & 30\\
Within-family commutators & 108\\
Class-two relations & 3732\\
Steinberg commuting relations & 12\\
Steinberg product relations & 12\\
Forward and inverse conjugation rules & 504\\
Negative-label names & 18\\
Distinguished-word name & 1\\
\midrule
Total & 4417\\
\bottomrule
\end{tabular}
\caption{The total freely reduced length is 41\,953 and the maximum is 19.
Free reduction removes adjacent inverse letters only; cyclic conjugates are not identified.}
\end{table}

\FloatBarrier
\subsection{Universal kernel and a fixed-error tolerance}
\begin{proposition}\label{prop:universal}
Every homomorphism $\pi:\Dhat\to\U(\M)$ sends $j$ to $I$. The same is true for the
nontrivial image of $j$ in $\Gamma$.
\end{proposition}
\begin{proof}
We follow Thom's proof of his Theorem~1.3~\cite[Section~5.2]{thom}.
Put $A=\pi(\Hhat)'\cap\M$. Both copies conjugate every $h'\in\Hhat$ by $t$ to the same
element of the shared subgroup. This holds for the twelve generators by
Eq.~\eqref{eq:conj-full}, and hence for all their words. Therefore
$c=t_0^{-1}t_1$ centralizes $\Hhat$, so $\pi(c)\in A$.

Proposition~\ref{prop:kazhdan-pair} verifies the hypotheses of
Theorem~\ref{thm:external} in copy 0. Consequently
\[
 \pi(t_1t_0^{-1})=\pi(t_0)\pi(c)\pi(t_0)^*\in A.
\]
It commutes with $\pi(h)$, giving $\pi(j)=I$. For a representation of $\Gamma$, compose
with $\Dhat\to\Gamma$.
\end{proof}

\begin{corollary}\label{cor:tolerance}
There is a rational $e_0\in(0,1]$ such that every finite-dimensional unitary assignment
$\varphi$ to the 43 positive generators satisfies
\begin{equation}
 \max_{r\in\Rel}\nnorm{\varphi(r)-I}_2\le e_0
 \quad\Longrightarrow\quad\nnorm{\varphi(j)-I}_2\le\frac12.
 \label{eq:e0}
\end{equation}
\end{corollary}
\begin{proof}
If no positive tolerance worked, for each $n$ there would be an assignment in some dimension
$m_n$ with relator defects at most $1/n$ and witness displacement greater than $1/2$.
The ultraproduct classes satisfy the presentation exactly and give a representation with
$\nnorm{\pi(j)-I}_2\ge1/2$, contradicting Proposition~\ref{prop:universal}. Once a positive
tolerance exists, choose a smaller positive rational one, at most one. This is the standard
compactness argument; compare the uniform statement in Ref.~\cite[Theorem~1.1]{liu}.
\end{proof}

The presentation and the witness threshold are fixed before the dimension is chosen.
No dimension cutoff occurs in Eq.~\eqref{eq:e0}.

\section{Exact word arithmetic and the original label table}\label{sec:arithmetic}

A Laurent polynomial is a finite list of exponent triples in $\Z^3$ with coefficients
in $\{1,2,3,4\}$. Combining like terms modulo five and removing zero coefficients gives
an exact representation. Integer action matrices remain over $\Z$; they are not reduced
modulo five. Besides Eq.~\eqref{eq:semidirect}, the inverse rule is
\begin{equation}
 (M,A)^{-1}=(\sigma_{A^{-1}}(M^{-1}),A^{-1}).
\end{equation}
A determinant-one matrix is inverted by its adjugate, so no general ring inversion procedure
is hidden in this rule.

Membership in $H_0$ is exactly the condition that the action matrix is $I$ and every
nonzero monomial in every matrix entry has nonnegative exponent coordinates. For sufficiency,
a polynomial determinant-one matrix has a polynomial adjugate, hence belongs to
$\SL_3(R_+)$. Using the full special linear group for $H_0$ avoids any need to identify
it with an elementary subgroup.

An amalgam word is a list of pairs consisting of a side and an element of $G_0$. Combine
adjacent syllables on the same side, remove identity syllables, and absorb a syllable of
$H_0$ into a neighbor. Each move shortens the list. A final alternating list of length at
least two with no syllable in $H_0$ is nontrivial; a one-syllable list is evaluated in $G_0$.
Equality of words $v,w$ is tested by reducing $v^{-1}w$. Literal equality of two reduced
lists is not used as a general equality criterion.

\subsection{Enumeration of labels and triangles}
Label 0 is the identity. Visit generators in the numbered order above,
each immediately followed by its inverse. Then visit relators in their stored order, and
their prefixes from left to right. Assign a new integer label at its first occurrence;
reuse an earlier label exactly when the quotient-word test identifies the two elements.
The named generator $j$ receives label 85. This label number is distinct from generator
number 43.

For a relator $s_1\cdots s_\ell$, include
\begin{equation}
 \bigl(\ell(s_1\cdots s_{k-1}),\ell(s_k),\ell(s_1\cdots s_k)\bigr)
 \qquad(2\le k\le\ell),
 \label{eq:prefixes}
\end{equation}
where $\ell(w)$ denotes a label, not word length. Also include
$(\ell(s),\ell(s^{-1}),0)$ for every positive generator. Keep the first occurrence of
each ordered triple; identity-input triples may be omitted, though none occur in this
instance. There are 37\,579 occurrences before deduplication, 16\,428 retained triples
in $\Tri_0$, and 11\,230 distinct labels.

\begin{lemma}\label{lem:prefix}
Unitary labels $w_g$ with $w_0=I$ satisfying every equation in $\Tri_0$ define a
homomorphism of $\Dhat$ by $s\mapsto w_{\ell(s)}$.
\end{lemma}
\begin{proof}
The generator-inverse equations give $w_{\ell(s^{-1})}=w_{\ell(s)}^{-1}$. Induction along
Eq.~\eqref{eq:prefixes} identifies a prefix label with the product of its assigned generator
matrices. The final prefix has identity label, so every relator holds. A removed duplicate
was already enforced; an identity-input equation is automatic after $w_0=I$.
\end{proof}

\section{Certified removal of 438 triangles}\label{sec:pruning}

The labels and the group presentation do not change in the new reduction. Only redundant
multiplication equations are removed. Protect the set
\[
 \mathcal F=\{(\ell(s),\ell(s^{-1}),0):s=1,\ldots,43\}.
\]
For a label $y$ of a signed generator, write $y^\vee$ for its inverse label. The original
table contains the protected equation in one of the orientations $(y,y^\vee,0)$ or
$(y^\vee,y,0)$.

For each pair of triples $(x,y,z)$ and $(z,y^\vee,x)$ in $\Tri_0$, retain the
lexicographically smaller triple and delete the larger. The program checks that no protected
triple occurs in these pairs. There are 408 such deletions. Next remove each zero-output
triple $(x,y,0)$ outside $\mathcal F$ whose reversed triple $(y,x,0)$ belongs to
$\mathcal F$; there are 30. All retained triples keep their original order.

\subsection{Free-group certificates}
Let $X_g$ be a free generator for label $g\ne0$ and put $X_0=1$. A triple $t=(x,y,z)$
has relation word
\[
 r_t=X_xX_yX_z^{-1}.
\]
For a retained relation $xy=z$ and inverse equation $q=X_yX_{y^\vee}$,
\begin{equation}
 r_t^{-1}\,X_xqX_x^{-1}=X_zX_{y^\vee}X_x^{-1}.
 \label{eq:delete-cert}
\end{equation}
This is an identity in the free group. If the protected inverse equation instead reads
$q'=X_{y^\vee}X_y$, use $q=X_yq'X_y^{-1}$. Thus the first family of deletions has
certificates with two conjugates of retained equations. A reversed inverse is certified by
\begin{equation}
 X_x(X_yX_x)X_x^{-1}=X_xX_y.
\end{equation}
The second family has one cell per certificate.

The stored certificate convention uses integer $g+1$ for $X_g$, with label 0 represented
by the empty word. Each cell has a retained original triangle index, a sign $\pm1$, and a
conjugating word. The verifier multiplies all cells, freely reduces the result, and compares
it with the relation word of the deleted triple. It also checks that no referenced equation
has itself been deleted.

\begin{proposition}\label{prop:pruning}
The 15\,990 retained equations imply all 16\,428 original equations. Every label occurs in
the retained table. Consequently Lemma~\ref{lem:prefix} still holds when only the retained
equations are assumed.
\end{proposition}
\begin{proof}
The 438 certificate identities prove the missing equations. The certificate list and retained
indices form a disjoint partition of $\Tri_0$, and every cell uses a retained equation.
Coverage is a finite check of the retained label indices. Applying Lemma~\ref{lem:prefix}
proves the final assertion.
\end{proof}

Coverage, rather than a claimed minimality property, is what the channel proof needs. No
assertion is made that this is the smallest sufficient triangle table.

\section{The channel family and exact extraction}\label{sec:channels}

Let $a=3/5$ and $b=4/5$. A reference list $B$ contains at least the identity label 0.
Put one entry $(A_g)_{pp}=1$ at each reference position $p$ labeled by $g$. A triangle
$(x,y,z)$ uses two new system positions $p,p+1$ with
\begin{equation}
 (A_0)_{p,p}=a,\qquad (A_y)_{p,p+1}=b,\qquad
 (A_x)_{p+1,p}=b,\qquad (A_z)_{p+1,p+1}=-a.
 \label{eq:matrix-entries}
\end{equation}
Different blocks have disjoint system positions; all unspecified entries are zero.
Set $\Phi(X)=\sum_gA_gXA_g^*$.

The original channel $\Phi_\full$ uses $\Tri_0$ and one singleton for every label.
The reduced channel $\Phi_1$ uses $\Tri$ with $B=(0)$; $\Phi_2$ uses $B=(0,85)$.
Their natural dimensions are $1+2\cdot15990=31981$ and $2+2\cdot15990=31982$.
For $\Phi_{15}$ append 786 singleton positions with label 0 to $\Phi_2$. These occupy
positions $31982,\ldots,32767$. The defining Kraus matrix $A_0$ receives entry one on
each added diagonal position. All other Kraus matrices vanish there.

\begin{proposition}\label{prop:channel}
Each channel is unital, trace preserving, and generally tracially factorizable, and has
Choi rank 11\,230. Its nonzero Choi eigenvalues are $b_g/(25d)$, with $b_g$ the sum of
squared integer Kraus numerators for label $g$.
\end{proposition}
\begin{proof}
Let $\lambda$ be the regular representation of $\Gamma$. The triangle block
\[
 \begin{pmatrix}aI&b\lambda(y)\\b\lambda(x)&-a\lambda(z)\end{pmatrix}
 =\operatorname{diag}(I,\lambda(x))
 \left(\begin{pmatrix}a&b\\b&-a\end{pmatrix}\otimes I\right)
 \operatorname{diag}(I,\lambda(y))
\]
is unitary, since $xy=z$ and $a^2+b^2=1$. Its direct sum with the reference unitaries is
$U=\sum_gA_g\otimes\lambda(g)$. The canonical trace satisfies
$\tau(\lambda(g)\lambda(h)^*)=\delta_{g,h}$, giving the stated Kraus channel and a general
tracial factorization by the criterion of Ref.~\cite[Theorem~2.2]{hm11}, up to their
convention of conjugating by $u^*$.

Every Kraus matrix has nonempty support, by coverage, and different labels have disjoint
supports. Hence the $A_g$ are linearly independent and Hilbert--Schmidt orthogonal. Therefore
\[
 J(\Phi)=d^{-1}\sum_g\vket{A_g}\vbra{A_g},\qquad
 \operatorname{spec}_{>0}J(\Phi)=\left\{\frac{b_g}{25d}:g\in\Lab\right\}.
\]
Each $A_g$ is a weighted partial permutation. In a triangle with nonidentity inputs,
$z=x$ would force $y=1$ and $z=y$ would force $x=1$, so neither occurs. The possible
coincidences $x=y$ and $z=0$ occupy different rows and columns within their Kraus matrix.
Different blocks are disjoint. Thus both normalization identities can also be checked from
integer row and column squared sums, each equal to 25.
\end{proof}

\subsection{Support determines the bath block array}
Let $\Pi=\supp J(\Phi)$. For a unitary realization $V$ in $M_d\otimes N$, write
$V_{pq}$ for its bath blocks. The Choi entry at $(p,q),(r,s)$ is
$d^{-1}\tau(V_{pq}V_{rs}^*)$.

\begin{lemma}\label{lem:support-extract}
If $\supp J(\Psi)\le\Pi$ for the channel implemented by $V$, then
$V=\sum_gA_g\otimes v_g$ for unitaries $v_g\in N$. The labels
$w_g=v_0^*v_g$ satisfy $w_0=I$ and $w_xw_y=w_z$ for every retained triple.
If $\Psi=\Phi$, then $\tau(v_gv_h^*)=\delta_{g,h}$.
\end{lemma}
\begin{proof}
The range of $\Pi$ is the span of the vectorized $A_g$. A scalar vector perpendicular to
that span has zero quadratic form in $J(\Psi)$, so its corresponding linear combination of
bath blocks has zero $2$-norm. Faithfulness makes that combination zero. A basis of these
linear constraints puts the full block array in the span of the $A_g$, giving
$V=\sum_gA_g\otimes v_g$.

The singleton for label 0 has no other nonzero entry in its row or column, hence $v_0$ is
unitary. Every triangle block is
\[
 \begin{pmatrix}av_0&bv_y\\bv_x&-av_z\end{pmatrix}.
\]
Its first row gives $a^2I+b^2v_yv_y^*=I$ and its first column gives
$a^2I+b^2v_x^*v_x=I$. A finite tracial algebra has no proper isometries: if $v^*v=I$,
then $I-vv^*$ is a projection of trace zero. Thus $v_x,v_y$ are unitary, and the second
column gives $v_z^*v_z=I$. All labels are covered. Row orthogonality gives
$v_0v_x^*=v_yv_z^*$, equivalently $v_z=v_xv_0^*v_y$, and normalization gives the
multiplication equation for $w$.

When the Choi states are equal, comparison in the orthogonal Kraus-vector basis gives
$\tau(v_gv_h^*)=\delta_{g,h}$.
\end{proof}

\section{Excluding every convergent finite-bath sequence}\label{sec:global}

\begin{lemma}\label{lem:ultra-factor}
If a channel $\Phi$ lies in $K_d$, it has an exact unitary realization in $M_d\otimes\M$
for a tracial matrix ultraproduct $\M$.
\end{lemma}
\begin{proof}
Choose finite matrix-tracial factorizations converging to $\Phi$. The system dimension $d$
is fixed. Every bath block is an operator-norm contraction, so its ultraproduct class exists.
The finitely many classes form a unitary because the original unitary identities are exact.
The trace of each bath-block product converges to its specified Choi entry, giving exactly
$\Phi$. Compare the implication (2)$\Rightarrow$(1) of Ref.~\cite[Theorem~3.6]{hm15}, which is
proved there through condition (3), with a factorization through the ultrapower of the
hyperfinite factor; here the matrix ultraproduct is taken directly.

For completeness, finite mixtures give no larger closure; see also
Ref.~\cite[Proposition~3.8(ii)]{musat}. For rational mixture weights
$n_i/N$ and bath dimensions $m_i$, choose a common multiple $L$ of the $m_i$ and use
$n_iL/m_i$ copies of the $i$th factorizing unitary. The total bath dimension is $NL$ and
its normalized trace gives weight $n_i/N$. Approximation of real weights completes the
reduction.
\end{proof}

\begin{theorem}\label{thm:nonapprox}
Each of $\Phi_\full$, $\Phi_1$, $\Phi_2$, and $\Phi_{15}$ lies in
$\Fact_d\setminus K_d$ and has $\Delta_J(\Phi)>0$.
\end{theorem}
\begin{proof}
Proposition~\ref{prop:channel} gives general factorization. Membership in $K_d$ would give
an exact ultraproduct factorization by Lemma~\ref{lem:ultra-factor}. Extract the unitary
labels using Lemma~\ref{lem:support-extract}. The pruning certificates and the original
prefix argument give a representation of $\Dhat$. Proposition~\ref{prop:universal} forces
$w_{85}=I$. But labels 0 and 85 are distinct and the target trace table gives
$\tau(w_{85})=\tau(v_0^*v_{85})=0$, contradicting $\tau(I)=1$.
Closedness of $K_d$ in the finite-dimensional space of channels gives positive distance.
\end{proof}

This argument does not assume that the finite approximating unitaries already have the
triangle block form. The exact target forces that form after passage to the ultraproduct.
In particular, excluding exact finite-dimensional representations alone would not prove the
theorem.

\subsection{A factor realization}
Let $F_2$ be the free group on two generators. The inclusion $\Gamma<\Gamma*F_2$ preserves the canonical trace, so the same
controlled unitary realizes every channel in $L(\Gamma*F_2)$. The free product is
countable and has infinite conjugacy classes: suitable reduced words in the two free
generators give infinitely many distinct conjugates of any nonidentity word. Thus its group
von Neumann algebra is a type $\mathrm{II}_1$ factor with separable predual.

If it embedded trace-preservingly into the hyperfinite ultrapower $R^\omega$, its finite
unitary tuples would have simultaneous matrix approximations to arbitrary finite lists of
moments. Coordinatewise approximation in finite-dimensional subalgebras of the hyperfinite
factor, followed by a diagonal choice, would give a matrix-ultraproduct representation of
$\Dhat$ retaining the trace $\tau(j)=0$. This contradicts Proposition~\ref{prop:universal}.
At the level of channels this is the implication (1)$\Rightarrow$(2) of
Ref.~\cite[Theorem~3.6]{hm15}, and at the level of groups the standard correspondence between
hyperlinear groups and embeddable group von Neumann algebras~\cite[Proposition~7.1]{ozawa}.
Liu's manuscript draws the corresponding non-embedding for the group von Neumann algebra of a
wreath product from its Theorem~1.1~\cite{liu}. By the same implication, no finite von Neumann
algebra that serves as an ancilla for one of these channels embeds trace-preservingly into
$R^\omega$; compare Ref.~\cite[Corollary~3.11]{musat}. The nonembeddable algebra is the bath factor, not the
finite system algebra $M_d$.

\section{The exact-support numerical bounds}\label{sec:restricted}

The following is the finite-dimensional form of Thom's observation that a compressor
normalizes every finite image of the subgroup~\cite[Remark~4.4]{thom}, combined with his
argument for the witness~\cite[Section~5.2]{thom}.

\begin{lemma}[Exact finite-dimensional rigidity]\label{lem:finite-rigid}
Let $H<G$ and $tHt^{-1}\le H$. In $D=G*_HG$, put $j=[h,t_1t_0^{-1}]$ for $h\in H$.
Every homomorphism $\pi:D\to\U(m)$ sends $j$ to $I$.
\end{lemma}
\begin{proof}
For $A=\pi(H)'\cap M_m$, compression gives
$\pi(t_0)A\pi(t_0)^*\supseteq A$. The two algebras have the same finite vector-space
dimension, so the inclusion is equality. Since $t_0^{-1}t_1$ centralizes $H$, its image
lies in $A$; conjugating by $\pi(t_0)$ puts $\pi(t_1t_0^{-1})$ in $A$. It commutes with
$\pi(h)$.
\end{proof}

The same proof works in a finite direct sum of matrix algebras. For our presentation, the
compression assumption follows directly from the twelve positive-root conjugation rules;
property $(T)$ is not needed. The dimension argument does not extend to a relative commutant
inside a matrix ultraproduct, which can be infinite-dimensional.

Let $T_d^\Pi=\{\Psi\in T_d:\supp J(\Psi)\le\Pi\}$, where $\Pi=\supp J(\Phi)$.
This class is nonempty: assigning every label the identity produces a unitary direct sum
of scalar blocks whose Choi support is contained in $\Pi$.

\begin{theorem}\label{thm:restricted}
Put $s_0=\nnorm{A_0}_{\mathrm{HS}}^2$ and $s_j=\nnorm{A_{85}}_{\mathrm{HS}}^2$. Then
\begin{equation}
 \inf_{\Psi\in\overline{T_d^\Pi}}\dJ(\Phi,\Psi)\ge\frac{\sqrt{s_0s_j}}{d}.
 \label{eq:restricted}
\end{equation}
If the two singleton references 0 and 85 are present, then also
\begin{equation}
 \inf_{\Psi\in\overline{T_d^\Pi}}\dd(\Phi,\Psi)\ge\frac12.
 \label{eq:restricted-half-diamond}
\end{equation}
\end{theorem}
\begin{proof}
Lemma~\ref{lem:support-extract} gives unitary labels for each $\Psi\in T_d^\Pi$.
The original or derived multiplication equations give the finite presentation, and
Lemma~\ref{lem:finite-rigid} gives $v_{85}=v_0$. Let
$\ket{a_0}=\vket{A_0}/\sqrt{s_0}$ and $\ket{a_j}=\vket{A_{85}}/\sqrt{s_j}$.
These vectors are orthogonal. Both Choi states have diagonal entries $s_0/d,s_j/d$ on
their span. The target cross entry is zero and the comparison cross entry is
$\sqrt{s_0s_j}/d$. The rank-one effect
\[
 E=\frac12(\ket{a_0}+\ket{a_j})(\bra{a_0}+\bra{a_j})
\]
has expectation difference $\sqrt{s_0s_j}/d$, proving Eq.~\eqref{eq:restricted}.

For the half-diamond assertion, take the two singleton positions $p_0,p_j$ and set
$\ket{+}=(\ket{p_0}+\ket{p_j})/\sqrt2$. The target output is
$(\ket{p_0}\bra{p_0}+\ket{p_j}\bra{p_j})/2$. Since $v_{85}=v_0$, the comparison
acts as the identity on this two-dimensional corner and outputs $\ket{+}\bra{+}$.
The difference has eigenvalues $1/2,-1/2$, so its half trace norm is $1/2$.
The projection $\ket{+}\bra{+}$ attains this expectation difference. No input-reference
entanglement is required. Continuity extends the two bounds to the stated closure.
\end{proof}

\begin{table}[ht]
\centering\small
\begin{tabular}{@{}lrrrr@{}}
\toprule
Instance & $d$ & $b_0$ & $b_{85}$ & Strict rational bound for $\dJ$\\
\midrule
Original $\Phi_\full$ & 44\,086 & 148\,534 & 41 & $1/447$\\
Reference-only reduction, two & 32\,858 & 148\,534 & 41 & $3/1000$\\
Reference-only reduction, one & 32\,857 & 148\,534 & 16 & $1/533$\\
Triangle-reduced $\Phi_1$ & 31\,981 & 144\,322 & 16 & $1/527$\\
Triangle-reduced $\Phi_2$ & 31\,982 & 144\,322 & 41 & $1/329$\\
Fifteen-qubit $\Phi_{15}$ & 32\,768 & 163\,972 & 41 & $1/316$\\
\bottomrule
\end{tabular}
\caption{All bounds compare with $\overline{T_d^\Pi}$, not with $K_d$.
The two reference-only variants are kept in the baseline part of the package. Each displayed bound is verified by integer arithmetic.}
\label{tab:gaps}
\end{table}

Indeed $s_g=b_g/25$, so a positive rational $q$ is certified by
$q^2<b_0b_{85}/(25d)^2$. For the fifteen-qubit instance the exact comparison is
\[
 316^2(163972\cdot41)>(25\cdot32768)^2.
\]
The complete comparison uses all finite bath dimensions but an \emph{exact} support
restriction. Arbitrarily close Choi states need not have exactly the same support.
Projecting onto $\Pi$ need not preserve a unitary finite-tracial realization. In particular,
neither $1/316$ nor $1/2$ is a global separation certificate.

\section{Global quantitative extraction in two metrics}\label{sec:quantitative}

Throughout this section the channel is fully referenced: each label has its own singleton
position $p_g$. Every relator and positive-generator inverse is encoded in the original
triangle table. Let $\ell_*\ge2$ be the maximum relator length. The statement below is
not asserted with unchanged constants for a reference-pruned or triangle-pruned table.

\begin{theorem}\label{thm:quantitative}
If $e_0\in(0,1]$ satisfies Eq.~\eqref{eq:e0}, then
\begin{equation}
 \Delta_J(\Phi_\full)\ge\frac{e_0^2}{4096\ell_*^2d},
 \qquad
 \Delta_\diamond(\Phi_\full)\ge\frac{e_0^2}{4096\ell_*^2}.
 \label{eq:quantitative-bounds}
\end{equation}
For $d=44086$ and $\ell_*=19$, the denominators are 65\,188\,028\,416 and 1\,478\,656,
respectively.
\end{theorem}

We prove both estimates, keeping their different input tests explicit. Let $U$ implement an
arbitrary finite matrix-tracial comparison $\Psi$ and put $B_g=U_{p_g,p_g}$.
For normalized Choi error $\delta=\dJ(\Phi_\full,\Psi)$ set $k=\sqrt{d\delta}$.
For half-diamond error $\eta=\dd(\Phi_\full,\Psi)$ set $k=\sqrt\eta$.
All bath blocks of $U$ are contractions.

\subsection{Singletons and individual entries}
In the Choi case, the projection onto entries outside the singleton's row diagonal has
zero target expectation. Its comparison expectation is at most $\delta$, yielding
\begin{equation}
 \tau(I-B_gB_g^*)\le d\delta=k^2.
 \label{eq:polar-choi}
\end{equation}
In the half-diamond case, input $\ket{p_g}\bra{p_g}$. The probability of leaving $p_g$
is at most $\eta$, so
\begin{equation}
 \tau(I-B_g^*B_g)\le\eta=k^2.
 \label{eq:polar-diamond}
\end{equation}
The two trace deficits agree. Complete the polar part of $B_g$ to a unitary $v_g$ in the
same matrix algebra. For singular values $0\le s\le1$, $(1-s)^2\le1-s^2$, hence
\begin{equation}
 \nnorm{B_g-v_g}_2\le k.
 \label{eq:polar-close}
\end{equation}

Take a triangle position $(a,b)$ carrying real coefficient $c\in\{3/5,4/5,-3/5\}$ and
label $g$. It is distinct from the singleton position. In the Choi case, the target
annihilates the scalar vector $\ket{a,b}-c\ket{p_g,p_g}$. Testing its normalized
rank-one projection gives
\[
 \nnorm{U_{ab}-cB_g}_2\le\sqrt{1+c^2}\,k.
\]
For the half-diamond case, input the normalized entangled state
\[
 \frac{\ket{b}\ket0+\ket{p_g}\ket1}{\sqrt2}.
\]
Its target output annihilates $\ket a\ket0-c\ket{p_g}\ket1$.
The expectation of the corresponding normalized projector is
$\nnorm{U_{ab}-cB_g}_2^2/[2(1+c^2)]$ for the comparison and at most $\eta$.
Thus
\[
 \nnorm{U_{ab}-cB_g}_2\le\sqrt{2(1+c^2)}\,k.
\]
Together with Eq.~\eqref{eq:polar-close}, both cases give
\begin{equation}
 \nnorm{U_{ab}-cv_g}_2<3k.
 \label{eq:entry-close}
\end{equation}
For the larger half-diamond constant this follows from
$\sqrt{2(1+c^2)}+|c|<3$, checked directly for $|c|=3/5,4/5$.

\subsection{Leakage in the row-orthogonality equation}
Let $S=\{p,p+1\}$ be the two system positions of a triangle. In the Choi case, a
projection onto positions in one of these rows and outside $S$ has zero target weight.
The sum of squared bath $2$-norms in that outside row is at most $d\delta=k^2$.
Since a block row of a unitary has operator norm at most one, the product of the two
outside rows has $2$-norm at most $k$.

In the half-diamond case, each input basis state in $S$ leaves $S$ with probability at
most $\eta$. Therefore
\[
 \sum_{a\notin S,\,b\in S}\nnorm{U_{ab}}_2^2\le2\eta.
\]
Unitarity of $U$ (each row and each column of bath blocks has total squared $2$-norm one)
gives equality of the total incoming and outgoing leakage sums:
\begin{align*}
 \sum_{a\in S,\,b\notin S}\nnorm{U_{ab}}_2^2
 &=|S|-\sum_{a,b\in S}\nnorm{U_{ab}}_2^2\\
 &=\sum_{a\notin S,\,b\in S}\nnorm{U_{ab}}_2^2.
\end{align*}
Each outside row consequently has squared $2$-norm at most $2\eta$. The product of
outside rows has $2$-norm at most $\sqrt2\,k$. Rectangular Hilbert--Schmidt norms here
use the original bath dimension as denominator.

The full row product is zero. Replacing the four within-block entries by their scalar
multiples of $v_g$, using Eq.~\eqref{eq:entry-close}, costs at most $12k$.
With $a=3/5,b=4/5$, this gives
\[
 ab\,\nnorm{v_0v_x^*-v_yv_z^*}_2
 \le\begin{cases}13k,&\text{Choi},\\(12+\sqrt2)k,&\text{half-diamond}.
 \end{cases}
\]
Since $ab=12/25$, both cases imply, for $w_g=v_0^*v_g$,
\begin{equation}
 \nnorm{w_xw_y-w_z}_2<32k.
 \label{eq:triangle-error}
\end{equation}
For example, replacing $\sqrt2$ by $3/2$ already makes the larger coefficient smaller
than 32.

\subsection{Relator error and the witness trace}
A relator of length $\ell$ has $\ell-1$ prefix merges. Equation~\eqref{eq:triangle-error}
therefore contributes at most $32(\ell-1)k$. A negative generator label differs from
the inverse of its positive label by at most $32k$, using its protected inverse equation.
Replacing every negative letter and setting $\varphi(s)=w_{\ell(s)}$ for positive
generators yields
\begin{equation}
 \max_{r\in\Rel}\nnorm{\varphi(r)-I}_2\le64\ell_* k.
 \label{eq:relator-error}
\end{equation}
The unitary $\varphi(j)$ is exactly the label $w_{85}$ because $j$ is a named positive
generator of the presentation.

The off-diagonal target entry between the identity and witness singletons is zero.
A traceless Hermitian matrix of trace norm $2\delta$ has operator norm at most $\delta$.
In the Choi case, its entry bound and Eq.~\eqref{eq:polar-close} give
\begin{equation}
 |\tau(w_{85})|\le d\delta+2\sqrt{d\delta}=k^2+2k.
 \label{eq:trace-choi}
\end{equation}
For the half-diamond case, use the input-reference state
$(\ket{p_0}\ket0+\ket{p_j}\ket1)/\sqrt2$. Its comparison cross entry is
$\tau(B_0B_{85}^*)/2$, so the same operator-norm observation gives
\begin{equation}
 |\tau(w_{85})|\le2\eta+2\sqrt\eta=2k^2+2k.
 \label{eq:trace-diamond}
\end{equation}
If the relevant distance is no larger than the right side of
Eq.~\eqref{eq:quantitative-bounds}, then $64\ell_*k\le e_0$ and $k\le1/4$.
But Eqs.~\eqref{eq:trace-choi}--\eqref{eq:trace-diamond} imply, respectively,
\[
 \nnorm{\varphi(j)-I}_2^2=2-2\operatorname{Re}\tau(w_{85})
 \ge\frac78\quad\text{or}\quad\frac34,
\]
both larger than $1/4$, contradicting Eq.~\eqref{eq:e0}.
This excludes every finite matrix-tracial comparison at the stated distance. Taking infima
and closure proves Theorem~\ref{thm:quantitative}.

The half-diamond estimate removes the system-dimension loss because its probes use only
the relevant input positions. It does not by itself provide the same dimension-free bound
in Choi distance. The norm comparison $\dd\le d\,\dJ$ recovers exactly the
$1/d$ factor in the Choi assertion.

\section{Why the tolerance is not supplied by a coordinate gap}\label{sec:gap}

For $m\ge2$, define
\[
 z_m=\frac{m^2-1+2mi}{m^2+1}.
\]
Assign one of the positive generators, for example generator 22, the matrix
$\operatorname{diag}(1,z_m)$ and all other positive generators $I_2$.
Exact rational arithmetic gives
\begin{equation}
 |z_m|^2=1,\qquad |1-z_m|^2=\frac4{m^2+1}.
\end{equation}
Since every relator has length at most 19, telescoping its unitary factors gives defect
at most $38/\sqrt{m^2+1}$.

For finite $m$, the tuple's commutant is the diagonal algebra. On the off-diagonal matrix
unit $E_{12}$, the conjugation energy
\[
 \frac1{4\cdot43}\sum_s\nnorm{[U_s,E_{12}]}_2^2
\]
divided by $\nnorm{E_{12}}_2^2$ is exactly
\begin{equation}
 \kappa_m=\frac1{43(m^2+1)}\longrightarrow0.
\end{equation}
Thus arbitrarily small presentation defects can coexist with a closing coordinate
conjugation gap. The witness approaches the identity too, so the example does not challenge
Proposition~\ref{prop:universal}. It rules out inserting an exact Kazhdan constant into
this coordinate estimate without an additional lifting or repair argument.

The internality theorem of Ref.~\cite{alt} gives no rates for the steps used here: its Proposition~2.2
chooses the Markov averaging times by a diagonal argument, and the later tolerances tend to
zero without rates. The normalization theorem of Ref.~\cite{thom} is likewise qualitative.
Neither supplies a numerical tolerance for the 4\,417 relators here. A global numerical certificate
therefore still requires one of two additional results: an evaluated tolerance in
Eq.~\eqref{eq:e0}, uniform over all dimensions, or a quantitative support-leakage theorem
that handles arbitrary nearby finite-tracial channels. Tests up to a finite bath dimension
cannot replace either result.

\section{Causal purity and the sharper memory bound}\label{sec:causal-full}

We use the causal service model of Ref.~\cite[Definition~2.3]{causal}. A service retains $Q$ qubits in an
initial state $\beta_0$, with purity $P=Q-S(\beta_0)$. Between rounds it exports
$j_t$ qubits permanently and imports $j_t$ qubits in a fixed joint state $\alpha_t$
independent of everything else. Put
\[
 c_t=j_t-S(\alpha_t),\qquad C=\sum_tc_t,\qquad J=\sum_tj_t.
\]
The qubits within an imported register may be correlated. Every retained register counts,
and there is no free reset, postselection, or measurement record. Each output is released
before the next input arrives. Error $\epsilon$ is the complete-experiment half trace
distance maximized over adaptive observers with quantum references. Entropies are in bits;
$\ln$ denotes the natural logarithm.

\subsection{Small entropy production gives a nearby tracial channel}
\begin{lemma}[{Local tracialization, Ref.~\cite[Lemma~5.4]{causal}}]\label{lem:tracialization}
For a unitary $U$ on $\C^d\otimes\C^M$ and bath state $\rho$, let
\[
 \Psi(X)=\Tr_M[U(X\otimes\rho)U^*],\qquad
 T(\rho)=\Tr_d[U(\tau_d\otimes\rho)U^*].
\]
Then $a=S(T(\rho))-S(\rho)\ge0$, and there is a finite mixture $\Theta$ of finite
matrix-tracial channels with
\begin{equation}
 \dJ(\Psi,\Theta)\le\sqrt{2\ln2\,a}.
 \label{eq:local-tracialization}
\end{equation}
\end{lemma}
\begin{proof}
Diagonalize input and output bath states with decreasing eigenvalues $\lambda_i$ and
$\mu_j$. If $U_{ji}\in M_d$ are the corresponding bath-indexed blocks, then
$t_{ji}=d^{-1}\Tr(U_{ji}^*U_{ji})$ is doubly stochastic and
$\mu_j=\sum_it_{ji}\lambda_i$. For probability distributions
$p_{ji}=t_{ji}\lambda_i$ and $q_{ji}=t_{ji}\mu_j$, the relative entropy in bits is
$D(p\Vert q)=S(\mu)-S(\lambda)=a$. In particular $a\ge0$, and Pinsker's inequality gives
\[
 T_0:=\sum_{ji}t_{ji}|\lambda_i-\mu_j|\le\sqrt{2\ln2\,a}.
\]
Pairing the ordered spectra minimizes the absolute-difference transport cost, so
$\sum_i|\lambda_i-\mu_i|\le T_0$.

With $\lambda_{M+1}=0$ and $P_k$ the top-$k$ input bath projection, write
\[
 \rho=\sum_{k=1}^Mw_kP_k/k,\qquad w_k=k(\lambda_k-\lambda_{k+1}).
\]
Let $Q_k$ be the top-$k$ output bath projection. The compression
$A_k=(I\otimes Q_k)U(I\otimes P_k)$ is a contraction between spaces of equal dimension
$dk$. Complete its polar factor to a unitary $V_k$, so $A_k=V_kH_k$, and put
\[
 W_{k,\pm}=V_k(H_k\pm i\sqrt{I-H_k^2}).
\]
These are unitaries. Averaging their conjugation channels replaces $A_kXA_k^*$ by
$A_kXA_k^*+C_kXC_k^*$, where $C_k=V_k\sqrt{I-H_k^2}$.
For input bath $P_k/k$, the averaged channel and the original one share the part from
$A_k$; their remaining positive Choi matrices have equal trace
$e_k=k^{-1}\sum_{i\le k<j}t_{ji}$. Their half trace distance is at most $e_k$.

Mixing these finite-tracial channels with weights $w_k$ gives $\Theta$ with
\begin{align*}
 \dJ(\Psi,\Theta)&\le\sum_kw_ke_k
 =\sum_{i<j}t_{ji}(\lambda_i-\lambda_j)\\
 &=\frac12\sum_{ij}t_{ji}|\lambda_i-\lambda_j|\\
 &\le\frac12\left(T_0+\sum_j|\mu_j-\lambda_j|\right)\le T_0.
\end{align*}
The equality with the half sum uses the vanishing signed sum supplied by double
stochasticity. This proves Eq.~\eqref{eq:local-tracialization}.
\end{proof}

\subsection{The entropy ledger}
Feed fresh maximally entangled input-reference pairs in every round and retain all
outputs and references. Let $\rho_t$ be the incoming memory marginal and $U_t$ the
round unitary. Its entropy production is
\[
 a_t=S\bigl(\Tr_d[U_t(\tau_d\otimes\rho_t)U_t^*]\bigr)-S(\rho_t).
\]
The round's output-reference marginal is unchanged after release, so the complete-experiment
error places it within $\epsilon$ of $J(\Phi)$. Lemma~\ref{lem:tracialization} and the
triangle inequality imply
\[
 a_t\ge\frac{(\Delta_J(\Phi)-\epsilon)_+^2}{2\ln2}.
\]
Exporting $j_t$ qubits lowers memory entropy by at most $j_t$; importing the independent
register raises it by $j_t-c_t$. Thus
$S(\rho_{t+1})\ge S(\rho_t)+a_t-c_t$. The initial memory entropy is $Q-P$ and the final
entropy is at most $Q$. Telescoping gives $\sum_ta_t\le P+C$ and therefore
\begin{equation}
 P+C\ge\frac{n(\Delta_J(\Phi)-\epsilon)_+^2}{2\ln2}.
 \label{eq:purity-full}
\end{equation}
This proof allows memory correlations with earlier outputs; it uses only the incoming
marginal and the independence of each fresh probe. There is no restriction on $Q$ or $J$.

\subsection{The flat-probe purity budget}
We reproduce the sharper form from Ref.~\cite{configuration}, Proposition 8.7, to make its
use independent of access to that companion paper.

\begin{lemma}[{Surprisal variance, Ref.~\cite[Lemma~5.9]{causal}; see also Ref.~\cite[Lemma~8.8]{configuration}}]\label{lem:variance}
If a state $\beta$ on $q$ qubits has eigenvalues $\beta_i$ and deficit
$D=q-S(\beta)$, the random variable $-\log_2\beta_i$, sampled with probability $\beta_i$,
has variance at most $(q+2/\ln2)D$.
\end{lemma}
\begin{proof}
For $x\ge0$ put $\phi(x)=x\ln x-x+1$. Then
\[
 x(\ln x)^2\le(2+\ln_+x)\phi(x).
\]
To check this, set $x=e^t$. For $t\ge0$, the difference is
$e^t(t-2)+t+2$, which and whose first derivative vanish at zero and whose second derivative
is $te^t\ge0$. For $t\le0$, the difference is
$2(te^t-e^t+1)-t^2e^t$, whose derivative is $-t^2e^t\le0$ and whose value at zero is
zero. The limit $x=0$ is immediate.

Let $m=2^q$ and $x_i=m\beta_i$. Their average is one, $x_i\le m$, and
$m^{-1}\sum_i\phi(x_i)=D\ln2$. Variance is unchanged by the additive constant $q$ and
is at most the corresponding second moment. Therefore
\[
 \operatorname{Var}(-\log_2\beta_i)
 \le\frac1{m(\ln2)^2}\sum_i x_i(\ln x_i)^2
 \le\left(q+\frac2{\ln2}\right)D.
\]
\end{proof}

\Needspace{9\baselineskip}
\begin{proposition}[{Flat-probe budget, Ref.~\cite[Proposition~8.7]{configuration}}]\label{prop:budget}
Suppose a pure input-reference probe has output $\omega$ flat of rank $r$ under $\Phi$.
Every causal service of $n$ uses with $\epsilon\le1/16$ and $J\le n\log_2r/2$ satisfies
\begin{equation}
 P+C<3Q+6.\label{eq:budget}
\end{equation}
\end{proposition}
\begin{proof}
Put $b=P+C$ and suppose $b\ge3Q+6$. Preload all independent imports and keep all exports
on the device side of a mathematical dilation. The device then starts with $L=Q+J$ qubits
in the product state $\sigma_0=\beta_0\otimes\bigotimes_f\alpha_f$. This reorganization
does not change the observer's reduced state. Every factor has at most $Q$ qubits.

If $\sigma_i$ are the eigenvalues, define $Z=L+\log_2\sigma_i$, sampled with probability
$\sigma_i$. Then $\mathbb EZ=b$. Independence and Lemma~\ref{lem:variance} give
\[
 \operatorname{Var}Z\le\left(Q+\frac2{\ln2}\right)b.
\]
Feed $n$ copies of the given pure probe. Conditioned on an eigenvector of $\sigma_0$, the
observer's final state $\tau_i$ has rank at most $2^L$, since it is a marginal of a pure
state with an $L$-qubit device. For a real $k$, there are at most $2^{L-k}$ eigenvalues
$\sigma_i\ge2^{k-L}$. The support projection $E$ of the corresponding $\tau_i$ therefore
has rank at most $2^{2L-k}$ and
$\Tr(E\tau_O)\ge\Pr(Z\ge k)$. The ideal state $\omega^{\otimes n}$ is flat of rank
$r^n$. Hence, using $2J\le n\log_2r$,
\begin{equation}
 \Pr(Z\ge k)\le\epsilon+2^{2Q-k}.\label{eq:tail}
\end{equation}

Set $u=b-2Q-3\ge b/3+1>0$. The one-sided Chebyshev inequality gives
$\Pr(Z\ge b-u)\ge u^2/(\operatorname{Var}Z+u^2)$. Equation~\eqref{eq:tail}, at
$k=b-u=2Q+3$, bounds the same probability by $1/16+1/8=3/16$.
Thus $13u^2\le3\operatorname{Var}Z$ and
\[
 b<\frac{27}{13}\left(Q+\frac2{\ln2}\right)<3Q+6,
\]
a contradiction. The final constant follows from $\ln2>9/13$; for an exact check, the first
three positive terms of $2\sum_{k\ge0}(1/3)^{2k+1}/(2k+1)$ already exceed $9/13$.
\end{proof}

\subsection{Application to the original and compact channels}
On the 11\,230 singleton positions of $\Phi_\full$, the channel is complete dephasing.
Half of a maximally entangled state on these positions gives the output
\[
 \frac1r\sum_{g\in\Lab}\ket{p_gp_g}\bra{p_gp_g},\qquad r=11\,230.
\]
This is flat of rank $r$. Its complement has minimal dimension $r$, so the maximal
complementary entropy is $\log_2r$.

Combining Eqs.~\eqref{eq:purity-full} and \eqref{eq:budget} gives
\begin{equation}
 Q\ge\frac{n(\Delta_J(\Phi_\full)-\epsilon)_+^2}{6\ln2}-2
 \qquad\left(\epsilon\le\frac1{16},\ J\le\frac n2\log_2r\right).
 \label{eq:memory-exact}
\end{equation}
At $\epsilon\le\Delta_J(\Phi_\full)/2$, this becomes
\begin{equation}
 Q\ge\frac{n\Delta_J(\Phi_\full)^2}{24\ln2}-2.
\end{equation}
With the earlier budget $P+C\le14Q+2\log_2(16/7)$ of Ref.~\cite[Theorem~5.8]{causal} in place
of Proposition~\ref{prop:budget}, the same argument gives the leading coefficient
$1/(112\ln2)$; Proposition~\ref{prop:budget} makes it $14/3$ times larger, although the
additive constant is worse. The rank-rate restriction is not asserted to be the optimal
exchange rate for a channel outside $K_d$.

For $\Phi_2$ and $\Phi_{15}$, the equal superposition of the two reference positions is
a pure probe with flat rank-two output. Thus Eq.~\eqref{eq:memory-exact} holds for either
channel, with its own global Choi gap and the budget $J\le n/2$. This does not restore the
original rank-$11\,230$ probe. No result here identifies the maximal complementary entropy
of a reference-pruned channel.

The support-restricted bounds of Appendix~\ref{sec:restricted} and the half-diamond bound of
Appendix~\ref{sec:quantitative} cannot be inserted as $\Delta_J$ in
Eq.~\eqref{eq:purity-full} without the appropriate global Choi comparison.

\section{Verification package and data format}\label{sec:files}

Only Python 3.10 or later and its standard library are needed for the finite-data checks.
The LaTeX sources are separate from those executable requirements. All paths in this
section are relative to the \path{verification} folder of the repository. From that folder,
run
\begin{lstlisting}
python verify_release.py
\end{lstlisting}
This runs the two retained baseline suites, the independent reduced-data checker, scalar
and probe checks, deliberate-corruption controls in normal and optimized Python, and
deterministic re-emission. The run writes its record to
\path{reports/verification.json}. The baseline suites retain their own detailed records.
The quicker command \texttt{python verify\_release.py -{}-quick} omits those baseline suites
and is not the complete verification.

\subsection{Where the objects are stored}
\begin{longtable}{@{}p{0.47\textwidth}p{0.48\textwidth}@{}}
\toprule
File & Content\\
\midrule
\endhead
\path{baseline/verify/data/qi218/presentation.json} & The original literal presentation, generator names,
relator families, and distinguished word.\\[4pt]
\path{baseline/verify/data/qi219/channel.json.gz} & The original fully referenced channel, all label words
and exact normal forms, generator labels, and 16\,428 triples.\\[4pt]
\path{data/deletion_certificates.json} & The 15\,990 retained original indices and all 438 free-group
certificates for deleted equations.\\[4pt]
\path{data/one_anchor_channel.json.gz} & The reduced channel $\Phi_1$ on $M_{31981}$.\\[4pt]
\path{data/two_anchor_channel.json.gz} & The reduced channel $\Phi_2$ on $M_{31982}$.\\[4pt]
\path{data/fifteen_qubits_channel.json.gz} & The main channel $\Phi_{15}$ on $M_{32768}$, including
all 786 added identity positions.\\[4pt]
\path{data/summary.json} & Construction counts, canonical hashes, restricted-gap squares, and explicit
null fields for numerical global gaps.\\[4pt]
\path{reports/finite_instances.json} & Written by the full run: independently computed Choi histograms, checked dimensions,
and the rational restricted bounds of Table~\ref{tab:gaps}.\\[4pt]
\path{channel_api.py} & Exact action on sparse inputs and matrix units; no dense channel matrix is
allocated.\\
\bottomrule
\end{longtable}

All system indices are zero-based. An entry \path{[row,column,numerator]} in
\path{kraus[g]} means the entry $\mathrm{numerator}/5$ of $A_g$. The field
\path{generator_labels} is a list of pairs from signed generator numbers to label numbers.
The field \path{labels} retains the baseline word representatives and exact group values.
The triangle list consists of ordered triples of label numbers. For the padded channel,
\path{anchor_labels} lists the two original reference labels; the separate field
\path{identity_padding_count} records the 786 added identity singletons.

For any matrix unit, the channel action is exactly
\begin{equation}
 \Phi(E_{uv})=\frac1{25}\sum_g
 \sum_{(a,u,n)\in\mathrm{kraus}[g]}
 \sum_{(b,v,k)\in\mathrm{kraus}[g]}nk\,E_{ab}.
\end{equation}
For example,
\begin{lstlisting}
python channel_api.py --variant fifteen_qubits --matrix-unit 0 1
python channel_api.py --variant fifteen_qubits --matrix-unit 2 2
\end{lstlisting}
The first returns zero, the off-diagonal entry of the two-reference dephasing corner.
The second evaluates a nontrivial triangle column.

\subsection{Rebuilding and independent checks}
The baseline builder recreates the presentation and original channel from the finite rules.
The new emitter reads that exact table, produces the deletion certificates, and writes all
three reduced instances. To reconstruct the complete chain in a new directory, run
\begin{lstlisting}
python rebuild_all.py --out rebuilt
python verify_reduced.py --data rebuilt/data
\end{lstlisting}
The independent reduced-data checker imports no emitter code. It obtains the same matrices
by removing old system slots and reindexing the surviving entries; it then appends the
specified identity positions. It verifies every free-group certificate using its own stack
reduction, checks the protected inverse equations and label coverage, and recomputes both
normalization identities and the full Choi weight distribution.

Deliberate corruptions include a Kraus coefficient, a triangle, a certificate sign, a
certificate dependency on a deleted equation, a witness reference, a padding count, the
coefficient denominator, and a restricted bound relabeled as global. The runner requires the
corresponding semantic rejection, not merely any nonzero exit code. These tests also run
under \texttt{python -O}; none of the acceptance conditions is an optimizable assertion.

\subsection{Provenance and proof scope}
The baseline presentation has SHA-256
\begin{lstlisting}
c17a57463b1d3c327f13dab537b2a61e9cba300625d2009a90202022ac65adfe
\end{lstlisting}
and the uncompressed canonical baseline channel has SHA-256
\begin{lstlisting}
38150decf3d203e3eabbc5dee8d8a3f2d6390aa960c50f90fe4c87b804729b68
\end{lstlisting}
The corresponding hash of the fifteen-qubit channel is
\begin{lstlisting}
1236ae6e7d587c738baac7c0e5815c90ece46ff8c9392fd0012614aae6e1aa1b
\end{lstlisting}
A hash binds bytes; it is not a proof of an algebraic identity. Canonical uncompressed bytes
are compared because gzip encodings can vary across library versions.

The release includes written proofs, not a formalization in a proof assistant. Its finite
checks certify the concrete presentation data, the quotient computations, the sparse channel
identities, and the supplied free-group implications. They do not prove the external
centralizer theorem, the normalization theorem, or the general analytical estimates.
The analytical group and channel claims rest on the written proofs and the external theorems;
finite-data checks do not change their status.

The file \path{analytical_status.json} records each main claim and its dependency. Both
\path{global_numerical_gap} and \path{relator_tolerance_e0} are null.


\begin{thebibliography}{99}
\bibitem{stinespring}
W. F. Stinespring, Positive functions on $C^*$-algebras,
\emph{Proc. Amer. Math. Soc.} \textbf{6}, 211--216 (1955).
\bibitem{connes}
A. Connes, Classification of injective factors. Cases $\mathrm{II}_1$, $\mathrm{II}_\infty$,
$\mathrm{III}_\lambda$, $\lambda\ne1$,
\emph{Ann. Math.} \textbf{104}, 73--115 (1976).
\bibitem{hm11}
U. Haagerup and M. Musat, Factorization and dilation problems for completely positive maps
on von Neumann algebras, \emph{Commun. Math. Phys.} \textbf{303}, 555--594 (2011),
\href{https://arxiv.org/abs/1009.0778}{arXiv:1009.0778}.
\bibitem{hm15}
U. Haagerup and M. Musat, An asymptotic property of factorizable completely positive maps
and the Connes embedding problem, \emph{Commun. Math. Phys.} \textbf{338}, 721--752 (2015),
\href{https://arxiv.org/abs/1408.6476}{arXiv:1408.6476}.
\bibitem{mr21}
M. Musat and M. R\o rdam, Factorizable maps and traces on the universal free product of
matrix algebras, \emph{Int. Math. Res. Not. IMRN} \textbf{2021}(23), 17951--17970 (2021),
\href{https://arxiv.org/abs/1903.10182}{arXiv:1903.10182}.
\bibitem{musat}
M. Musat, The Connes--Kirchberg problem and infinite-dimensional phenomena in quantum
information theory, \href{https://arxiv.org/abs/2512.00432v1}{arXiv:2512.00432v1} (2025), to appear in
\emph{M\"unster J. Math.}
\bibitem{mip}
Z. Ji, A. Natarajan, T. Vidick, J. Wright, and H. Yuen,
$\mathrm{MIP}^*=\mathrm{RE}$, \href{https://arxiv.org/abs/2001.04383v3}{arXiv:2001.04383v3} (2022).
\bibitem{openai}
OpenAI, Nonsofic groups exist, Chapter~3 of \emph{Ten Advances in Mathematics and Theoretical
Computer Science} (2026), \url{https://cdn.openai.com/pdf/ten-proofs-oai.pdf}.
\bibitem{kt}
G. Kun and A. Thom, Nonsofic wreath products of residually finite groups,
\href{https://arxiv.org/abs/2608.06222v3}{arXiv:2608.06222v3} (2026), Theorems A, E and 4.1.
\bibitem{at}
V. Alekseev and A. Thom, Centralizers of sofic approximations of Kazhdan groups,
\href{https://arxiv.org/abs/2608.05362}{arXiv:2608.05362} (2026), Open Problem 6.2(a).
\bibitem{thom}
A. Thom, A conditional construction of a nonhyperlinear group and the centralizer problem,
\href{https://arxiv.org/abs/2609.27795v1}{arXiv:2609.27795v1} (14 September 2026), preprint,
Theorems 1.2 and 1.3, Remark 4.4, and Section 5.2.
\bibitem{liu}
J. Liu, Nonhyperlinear groups exist, manuscript dated 21 September 2026 (first version posted
20 September 2026), formerly at
\nolinkurl{https://jihaoliu.org/ai-results/nonhyperlinear-groups-exist-2026-09-20.pdf}, archived on
26 September 2026 at
\url{https://web.archive.org/web/20260926151036/https://jihaoliu.org/ai-results/nonhyperlinear-groups-exist-2026-09-20.pdf},
Theorems 1.1 and 1.2.
\bibitem{alt}
V. Alekseev, J. Liu, and A. Thom, Spectral gap, internality of relative commutants, and
non-hyperlinear groups, \href{https://arxiv.org/abs/2609.37761v1}{arXiv:2609.37761v1}
(29 September 2026), preprint, Theorem B, Corollary C, and Proposition 2.2.
\bibitem{higman}
G. Higman, Subgroups of finitely presented groups, \emph{Proc. Roy. Soc. London Ser. A}
\textbf{262}, 455--475 (1961).
\bibitem{ejz}
M. Ershov and A. Jaikin-Zapirain, Property $(T)$ for noncommutative universal lattices,
\emph{Invent. Math.} \textbf{179}, 303--347 (2010),
\href{https://arxiv.org/abs/0809.4095v2}{arXiv:0809.4095v2}, Theorems 1.2 and 6.2 and Corollary 4.7.
\bibitem{mr}
M. Musat and M. R\o rdam, Non-closure of quantum correlation matrices and factorizable
channels that require infinite dimensional ancilla (with an appendix by N. Ozawa),
\emph{Commun. Math. Phys.} \textbf{375},
1761--1776 (2020), \href{https://arxiv.org/abs/1806.10242}{arXiv:1806.10242}.
\bibitem{mori}
M. Mori, On the shape of correlation matrices for unitaries, \emph{Math. Scand.}
\textbf{130}, 359--363 (2024), \href{https://doi.org/10.7146/math.scand.a-142800}{doi:10.7146/math.scand.a-142800},
\href{https://arxiv.org/abs/2308.03345}{arXiv:2308.03345}.
\bibitem{ad}
C. Anantharaman-Delaroche, On ergodic theorems for free group actions on noncommutative spaces,
\emph{Probab. Theory Related Fields} \textbf{135}, 520--546 (2006).
\bibitem{pr}
V. I. Paulsen and M. Rahaman, Bisynchronous games and factorizable maps,
\emph{Ann. Henri Poincar\'e} \textbf{22}, 593--614 (2021),
\href{https://arxiv.org/abs/1908.03842}{arXiv:1908.03842}.
\bibitem{bhv}
B. Bekka, P. de la Harpe, and A. Valette, \emph{Kazhdan's Property (T)}, New Mathematical
Monographs \textbf{11} (Cambridge University Press, 2008).
\bibitem{slofstra}
W. Slofstra, Tsirelson's problem and an embedding theorem for groups arising from non-local
games, \emph{J. Amer. Math. Soc.} \textbf{33}, 1--56 (2020),
\href{https://arxiv.org/abs/1606.03140}{arXiv:1606.03140}.
\bibitem{ozawa}
N. Ozawa, About the QWEP conjecture, \emph{Internat. J. Math.} \textbf{15}, 501--530 (2004),
\href{https://arxiv.org/abs/math/0306067}{arXiv:math/0306067}.
\bibitem{causal}
S. Douglas and N. Mghirbi, Causal quantum-channel simulation: memory beyond entropy,
version 1.1.2 (2026), \href{https://doi.org/10.5281/zenodo.23070870}{doi:10.5281/zenodo.23070870},
Definition 2.3, Theorem 3.2, Lemmas 5.4 and 5.9, and Theorem 5.8.
\bibitem{configuration}
S. Douglas and N. Mghirbi, Amenable groups with nearly exponential sofic profile, and quantum
channels that need nearly linear memory, version 1.0.2 (2026),
\href{https://doi.org/10.5281/zenodo.23070871}{doi:10.5281/zenodo.23070871},
Proposition 8.7 and Lemma 8.8.
\end{thebibliography}

\end{document}
